% ============================================================================
%  Turbulance: A Formal Specification of a Declarative Language for
%  Evidence-Bearing, Uncertainty-Aware Computation
%
%  Grammar, Type System, and Operational Semantics.
%
%  This document is self-contained: it defines the language in full from
%  first principles and cites only the standard external literature.  No
%  companion manuscript is required to understand the language.
% ============================================================================
\documentclass[11pt,a4paper]{article}

\usepackage[utf8]{inputenc}
\usepackage[T1]{fontenc}
\usepackage{lmodern}
\usepackage[a4paper,margin=2.4cm]{geometry}
\usepackage{amsmath,amssymb,amsthm,mathtools}
\usepackage{stmaryrd}      % \llbracket \rrbracket
\usepackage{enumitem}
\usepackage{booktabs}
\usepackage{array}
\usepackage{xcolor}
\usepackage{listings}
\usepackage{microtype}
\usepackage[numbers,sort&compress]{natbib}
\usepackage[colorlinks=true,linkcolor=blue!45!black,citecolor=blue!45!black,
            urlcolor=blue!55!black]{hyperref}
\usepackage{cleveref}

% ---------------------------------------------------------------------------
%  Theorem environments
% ---------------------------------------------------------------------------
\theoremstyle{plain}
\newtheorem{theorem}{Theorem}[section]
\newtheorem{lemma}[theorem]{Lemma}
\newtheorem{proposition}[theorem]{Proposition}
\newtheorem{corollary}[theorem]{Corollary}
\theoremstyle{definition}
\newtheorem{definition}[theorem]{Definition}
\newtheorem{convention}[theorem]{Convention}
\newtheorem{example}[theorem]{Example}
\theoremstyle{remark}
\newtheorem{remark}[theorem]{Remark}

% ---------------------------------------------------------------------------
%  Grammar / semantics notation
% ---------------------------------------------------------------------------
\newcommand{\nt}[1]{\ensuremath{\langle\textsf{#1}\rangle}}  % nonterminal
\newcommand{\tm}[1]{\textbf{\texttt{#1}}}                 % terminal keyword
\newcommand{\opt}[1]{\ensuremath{[\,#1\,]}}               % optional
\newcommand{\rep}[1]{\ensuremath{\{\,#1\,\}}}             % repetition (0+)
\newcommand{\bnfeq}{\;::=\;}
\newcommand{\bnfor}{\;\mid\;}
\newcommand{\rn}[1]{\textsc{#1}}                          % rule name
\newcommand{\sem}[1]{\ensuremath{\llbracket #1 \rrbracket}} % semantic brackets
\DeclareMathOperator{\dom}{dom}                           % domain of a map
\newcommand{\Downarrowx}{\Downarrow}                      % big-step
\newcommand{\evalto}{\Downarrow}
\newcommand{\stepto}{\longrightarrow}
\newcommand{\env}{\sigma}
\newcommand{\Val}{\mathbb{V}}
\newcommand{\Conf}{\mathcal{C}}
\newcommand{\Unit}{[0,1]}
\newcommand{\nor}{\oplus}                                 % noisy-OR
\newcommand{\att}{\ominus}                                % attenuation
\newcommand{\Type}{\tau}
\newcommand{\dyn}{\star}                                  % dynamic type
\newcommand{\subtype}{<:}
\newcommand{\consistent}{\sim}
\newcommand{\Gam}{\Gamma}
\newcommand{\nrm}{\mathsf{nrm}}
\newcommand{\ret}{\mathsf{ret}}
\newcommand{\wrong}{\mathbf{wrong}}
\newcommand{\Strat}{\varsigma}

% Inference-rule macro (portable; no external package required)
\newcommand{\infrule}[3]{%
  \ensuremath{\displaystyle\frac{\,#2\,}{\,#3\,}}\;\rn{(#1)}}
\newcommand{\jspace}{\qquad}

% ---------------------------------------------------------------------------
%  Listings: Turbulance
% ---------------------------------------------------------------------------
\definecolor{kwcol}{rgb}{0.20,0.20,0.65}
\definecolor{strcol}{rgb}{0.55,0.30,0.10}
\definecolor{comcol}{rgb}{0.35,0.45,0.35}
\lstdefinelanguage{turbulance}{
  morekeywords={funxn,item,proposition,motion,evidence,pattern,support,
    contradict,inconclusive,within,given,considering,for,each,in,while,
    return,ensure,allow,research,cause,point,resolution,resolve,cycle,drift,
    flow,roll,until,settled,over,on,goal,metacognitive,import,from,as,try,
    catch,finally,project,sources,otherwise,and,or,not,true,false,with_confidence,
    requirements,track,evaluate,derive_hypotheses,meta},
  sensitive=true,
  morecomment=[l]{//},
  morecomment=[l]{\#},
  morestring=[b]",
}
\lstset{
  language=turbulance,
  basicstyle=\ttfamily\small,
  keywordstyle=\color{kwcol}\bfseries,
  stringstyle=\color{strcol},
  commentstyle=\color{comcol}\itshape,
  showstringspaces=false,
  columns=fullflexible,
  keepspaces=true,
  breaklines=true,
  frame=single,
  framesep=4pt,
  rulecolor=\color{black!25},
  xleftmargin=10pt,
}

% ---------------------------------------------------------------------------
\title{\textbf{Turbulance: A Formal Specification of a Declarative Language\\
for Evidence-Bearing, Uncertainty-Aware Computation}\\[4pt]
\large Lexical Structure, Grammar, Type System, and Operational Semantics}

\author{Kundai Farai Sachikonye\\[2pt]
\small Technical University of Munich\\
\small \texttt{kundai.sachikonye@wzw.tum.de}}

\date{\today}

% ===========================================================================
\begin{document}
\maketitle

\begin{abstract}
\noindent
We give a complete, self-contained formal specification of \emph{Turbulance},
a declarative programming language whose primitive notions are not values and
side effects but \emph{evidence}, \emph{hypotheses}, and \emph{graded
uncertainty}. Turbulance unifies three computational paradigms that are
usually kept apart. From the imperative and functional traditions it inherits
bindings, first-class functions, and structured control flow. From logic and
argumentation it inherits \emph{propositions} carrying named \emph{motions}
that are confirmed or refuted by \emph{support} and \emph{contradict}
statements over observed data. From probabilistic and fuzzy computation it
inherits \emph{points}---semantic units annotated with a confidence and a
distribution over interpretations---together with \emph{resolution} operators
that collapse a point to a determinate value under an explicit strategy. We
formalise the language in five layers. (i) A lexical grammar over a fixed
token alphabet. (ii) A context-free \emph{concrete grammar} in extended
Backus--Naur form, with a fully specified operator-precedence hierarchy. (iii)
An \emph{abstract syntax} together with a desugaring translation. (iv) A
\emph{type system} that is gradual: it admits a dynamic type and a consistency
relation so that annotated and unannotated code interoperate, while
guaranteeing soundness on the fully annotated fragment. (v) An \emph{operational
semantics}, given in big-step form, comprising the deterministic functional
core, a confidence algebra on the unit interval, an argumentation semantics for
propositions, a resolution semantics for points, and rules for the four
``hybrid'' iteration constructs (\texttt{cycle}, \texttt{drift}, \texttt{flow},
\texttt{roll}). We then formalise the \emph{four-file project model}, in which a
program is a linked tuple of an orchestration unit, a dependency graph, a state
surface, and a decision log. Finally we establish the metatheory: determinacy
of evaluation, type safety (progress and preservation) for the static fragment,
boundedness and monotonicity of the evidence semantics, and termination
criteria for the iteration constructs. Throughout, the design is related to the
standard literature on formal languages, type systems, operational semantics,
probabilistic and fuzzy computation, and abstract argumentation. The
specification is intended to be precise enough to serve as an implementation
blueprint and stable enough to serve as a language standard.

\medskip
\noindent\textbf{Keywords:} programming-language specification, operational
semantics, gradual typing, probabilistic programming, fuzzy logic, abstract
argumentation, evidence theory, declarative languages, domain-specific
languages.
\end{abstract}

\tableofcontents
\bigskip
\hrule

% ===========================================================================
\section{Introduction}
\label{sec:intro}
% ===========================================================================

\subsection{The problem: computation that reasons under evidence}

Mainstream programming languages model computation as the transformation of
\emph{determinate} values: an expression denotes a value, a statement updates a
store, and a function maps inputs to outputs \citep{landin1966,winskel1993,
pierce2002}. This model is an excellent fit for systems whose inputs are known
and whose correctness criteria are crisp. It is a poor fit for a large and
growing class of tasks---scientific data analysis, evidence synthesis,
diagnostic reasoning---whose inputs are \emph{observations} of uncertain
provenance and whose outputs are \emph{conclusions} that are graded,
defeasible, and accountable to the evidence that produced them. In these tasks
the programmer is forced to encode, by hand and without language support, three
notions that the host language does not name: a \emph{hypothesis} under test, a
body of \emph{evidence} bearing on it, and a \emph{degree of belief} attached
to every intermediate result. The encoding is repetitive, error-prone, and
opaque; uncertainty is threaded through ordinary floating-point variables, and
the chain of reasoning that justifies a conclusion is lost.

Turbulance is a language in which these three notions are primitive. A
\emph{proposition} is a first-class declaration of a hypothesis, decomposed
into named \emph{motions}; a motion is confirmed or refuted by \emph{support}
and \emph{contradict} statements evaluated against data; and the confirmation
of a motion is not a Boolean but a confidence drawn from the unit interval,
accumulated by an explicit and inspectable argumentation procedure
\citep{toulmin1958,dung1995,shafer1976}. Orthogonally, a \emph{point} is a unit
of meaning---typically a fragment of text or a structured datum---that carries
both a confidence and a distribution over candidate \emph{interpretations}; a
\emph{resolution} operator collapses a point to a determinate value under an
explicit \emph{strategy} (maximum likelihood, Bayesian weighting, conservative
or exploratory selection, full distribution), so that the act of disambiguation
is itself a named, auditable computational step \citep{zadeh1965,kozen1981,
goodman2008}.

\subsection{Design principles}
\label{sec:principles}

The language obeys five design principles, each of which has direct
consequences for the formal definition that follows.

\begin{enumerate}[leftmargin=1.4em]
\item \textbf{Uncertainty is first-class, not encoded.} Confidence is a
semantic domain of the language (\cref{sec:domains}), not a convention on
\texttt{Float} values. Every point, motion, and resolution carries a value in
$\Unit$ with a fixed algebra of combination (\cref{sec:confalg}).

\item \textbf{Hypotheses are programs.} A proposition is a declaration with
operational meaning: evaluating it runs an argumentation procedure that returns
a verdict and a confidence (\cref{sec:propsem}). This realises, at the level of
language semantics, the view of inference as a defeasible debate rather than a
monotone deduction \citep{dung1995,besnard2008}.

\item \textbf{Determinism and probability cohere in one control flow.} The same
program text mixes deterministic iteration (\texttt{for}, \texttt{while}) with
probabilistic iteration (\texttt{considering}, \texttt{drift}) and convergent
iteration (\texttt{roll until settled}). The semantics treats these uniformly
as instances of a single iteration schema parameterised by how the loop
variable and its auxiliary confidence are bound (\cref{sec:hybrid}).

\item \textbf{Types are gradual.} Turbulance is dynamically typed with optional
annotations and strong inference. We model this with a dynamic type and a
consistency relation in the manner of \citet{siek2006}, obtaining
interoperability between annotated and unannotated code together with soundness
on the annotated fragment (\cref{sec:types,sec:meta}).

\item \textbf{A program is a four-part project.} Source is organised into four
linked file roles---an orchestration unit (\texttt{.trb}), a dependency graph
(\texttt{.ghd}), a state surface (\texttt{.fs}), and a decision log
(\texttt{.hre}). We formalise this as a typed module composition with a
linkage relation (\cref{sec:fourfile}).
\end{enumerate}

\subsection{Contributions and structure}

This document is the language standard. Its contributions are:

\begin{itemize}[leftmargin=1.4em]
\item a precise lexical grammar (\cref{sec:lexical}) and context-free concrete
grammar with a complete precedence hierarchy (\cref{sec:concrete});
\item an abstract syntax and a desugaring translation that reduces the surface
language to a small core (\cref{sec:abstract});
\item the semantic domains, including a confidence algebra on $\Unit$ and the
structures of points, floors, and corpora (\cref{sec:domains});
\item a gradual type system with judgements for the functional core, the
scientific-reasoning layer, and the probabilistic layer (\cref{sec:types});
\item a big-step operational semantics for all of the above, including the
argumentation semantics of propositions, the resolution semantics of points,
and the four hybrid iteration constructs (\cref{sec:dynamics});
\item the four-file execution model (\cref{sec:fourfile});
\item and the metatheory: determinacy, type safety, evidence-semantics
boundedness and monotonicity, and termination (\cref{sec:meta}).
\end{itemize}

\Cref{sec:examples} gives worked examples; \cref{sec:related} situates the
language with respect to existing paradigms; \cref{sec:conclusion} concludes.

\subsection{Metalanguage and notational conventions}
\label{sec:notation}

We use standard mathematical metalanguage. $\mathbb{N}$, $\mathbb{Z}$,
$\mathbb{R}$ are the naturals, integers, reals; $\Unit$ is the closed unit
interval. For a set $A$, $A^{*}$ is the set of finite sequences over $A$, and
$\mathcal{M}_{\mathrm{f}}(A)$ the set of finite multisets over $A$. A finite
partial map $f\colon A \rightharpoonup B$ has domain $\dom(f)$; $f[a\mapsto b]$
is the update; $\varnothing$ is the empty map. We write $\overline{x}$ for a
finite sequence $x_1,\dots,x_n$ whose length $|\overline{x}|=n$ is left
implicit. Grammars are written in extended Backus--Naur form
\citep{backus1959,naur1963}: $\nt{n}$ is a nonterminal, $\tm{k}$ a terminal,
$\opt{\alpha}$ an optional phrase, $\rep{\alpha}$ zero-or-more repetition,
$\alpha\bnfor\beta$ alternation. Inference rules are written
$\frac{\text{premises}}{\text{conclusion}}$ with a rule name; double bars are
not used. Semantic (denotation) brackets are $\sem{\cdot}$. Big-step evaluation
is written with $\evalto$.

\begin{convention}[Faithfulness]
The grammar, types, and semantics in this document describe the language as a
mathematical object. Where an implementation is mentioned it is only to fix the
intended reading; the present text, not any implementation, is normative.
\end{convention}

% ===========================================================================
\section{Lexical Structure}
\label{sec:lexical}
% ===========================================================================

A Turbulance program is a finite sequence of Unicode code points. Lexical
analysis partitions this sequence into a stream of \emph{tokens} after
discarding layout and comments. The lexical grammar is regular
\citep{kleene1956,thompson1968,hopcroft2006} and is given below as a set of
regular definitions; tokenisation follows the maximal-munch (longest-match)
rule, with keywords taking priority over identifiers.

\subsection{Layout and comments}

\emph{Whitespace} comprises the space, horizontal tab, carriage return, and
line feed. Whitespace separates tokens and is otherwise insignificant except
where noted (Turbulance uses explicit block delimiters; see
\cref{rem:layout}). A \emph{comment} begins with \texttt{//} or \texttt{\#}
and extends to the end of the line. Comments are lexically equivalent to
whitespace.

\subsection{Token classes}

\begin{definition}[Lexical alphabet]
The terminal alphabet $\Sigma_{\mathrm{lex}}$ consists of the token classes:
\emph{keywords}, \emph{identifiers}, \emph{numeric literals}, \emph{string
literals}, \emph{boolean literals}, \emph{operators}, and \emph{delimiters}.
\end{definition}

\paragraph{Identifiers and keywords.}
\begin{align*}
\nt{letter}     &\bnfeq \texttt{A}\bnfor\cdots\bnfor\texttt{Z}\bnfor
                        \texttt{a}\bnfor\cdots\bnfor\texttt{z}\bnfor\texttt{\_}\\
\nt{digit}      &\bnfeq \texttt{0}\bnfor\cdots\bnfor\texttt{9}\\
\nt{ident}      &\bnfeq \nt{letter}\,\rep{\nt{letter}\bnfor\nt{digit}}
\end{align*}
An \nt{ident} that coincides with a keyword is lexed as that keyword. The
\emph{keywords} of Turbulance are the reserved words
\[
\begin{array}{l}
\tm{funxn}\ \tm{item}\ \tm{proposition}\ \tm{motion}\ \tm{evidence}\
\tm{pattern}\ \tm{support}\ \tm{contradict}\ \tm{inconclusive}\\
\tm{within}\ \tm{given}\ \tm{considering}\ \tm{for}\ \tm{each}\ \tm{in}\
\tm{while}\ \tm{return}\ \tm{ensure}\ \tm{allow}\ \tm{research}\ \tm{cause}\\
\tm{point}\ \tm{resolution}\ \tm{resolve}\ \tm{cycle}\ \tm{drift}\ \tm{flow}\
\tm{roll}\ \tm{until}\ \tm{settled}\ \tm{over}\ \tm{on}\ \tm{goal}\\
\tm{metacognitive}\ \tm{import}\ \tm{from}\ \tm{as}\ \tm{try}\ \tm{catch}\
\tm{finally}\ \tm{project}\ \tm{sources}\ \tm{otherwise}\
\tm{requirements}\\
\tm{with\_confidence}\ \tm{true}\ \tm{false}\ \tm{and}\ \tm{or}\ \tm{not}.
\end{array}
\]
The keyword for function definition is \tm{funxn} (not \texttt{function}); the
keyword for binding is \tm{item} (not \texttt{var} or \texttt{let}); the
conditional keyword is \tm{given} (not \texttt{if}). These choices are
intrinsic to the language's vocabulary and are not abbreviations.

\paragraph{Literals.}
\begin{align*}
\nt{num}    &\bnfeq \nt{digit}\,\rep{\nt{digit}}\;
                    \opt{\texttt{.}\,\nt{digit}\,\rep{\nt{digit}}}\\
\nt{string} &\bnfeq \texttt{"}\;\rep{\nt{schar}}\;\texttt{"}\\
\nt{schar}  &\bnfeq \text{any code point except \texttt{"} and newline}
                    \bnfor \texttt{\textbackslash}\,\nt{esc}\\
\nt{esc}    &\bnfeq \texttt{n}\bnfor\texttt{t}\bnfor\texttt{r}\bnfor
                    \texttt{"}\bnfor\texttt{\textbackslash}\\
\nt{bool}   &\bnfeq \tm{true}\bnfor\tm{false}
\end{align*}
Numeric literals denote elements of $\mathbb{R}$ (the language does not
distinguish a separate integer type at the value level; see
\cref{sec:domains}); the type system nonetheless tracks a refinement
\textsf{Num} that an implementation may represent with machine integers when
the fractional part is absent.

\paragraph{Operators and delimiters.}
The operator tokens, in decreasing binding strength, are
\[
\texttt{.}\quad \texttt{!}\ \ (\text{unary})\quad \texttt{*}\ \texttt{/}\quad
\texttt{+}\ \texttt{-}\quad \texttt{<}\ \texttt{>}\ \texttt{<=}\ \texttt{>=}
\quad \texttt{==}\ \texttt{!=}\quad \texttt{\&\&}\quad \texttt{||}\quad
\texttt{|}\ \texttt{|>}\quad \texttt{=>}\quad \texttt{=}.
\]
The delimiters are \texttt{( ) \{ \} [ ]} and the punctuators
\texttt{,\ :\ ;\ .}\ . The colon \texttt{:} introduces a block (see
\cref{sec:concrete}); the semicolon is an optional statement terminator.

\begin{remark}[Block structure]
\label{rem:layout}
Turbulance delimits blocks explicitly. A block is either a brace-delimited
sequence \texttt{\{}\,$\dots$\,\texttt{\}} or a colon-introduced suite
\texttt{:}\,$\dots$ terminated by a dedent or an explicit block end, at the
discretion of the concrete syntax in \cref{sec:concrete}. To keep the grammar
context-free and layout-insensitive, we take the brace form as primitive and
treat the colon-suite form as syntactic sugar that an implementation's reader
expands to braces. All semantics in this document are stated over the
brace-delimited abstract syntax.
\end{remark}

% ===========================================================================
\section{Concrete Grammar}
\label{sec:concrete}
% ===========================================================================

We now give the context-free grammar of Turbulance \citep{chomsky1956,aho2006}.
The grammar is presented in EBNF; \cref{sec:abstract} maps it to abstract
syntax. The start symbol is \nt{program}.

\subsection{Programs and declarations}

\begin{align*}
\nt{program} &\bnfeq \rep{\nt{decl}}\\[2pt]
\nt{decl}    &\bnfeq \nt{funDecl} \bnfor \nt{propDecl} \bnfor \nt{evidenceDecl}
                     \bnfor \nt{patternDecl} \bnfor \nt{pointDecl}\\
             &\bnfor \nt{goalDecl} \bnfor \nt{metaDecl} \bnfor \nt{importDecl}
                     \bnfor \nt{projectDecl} \bnfor \nt{stmt}\\[4pt]
\nt{funDecl} &\bnfeq \tm{funxn}\ \nt{ident}\
                     \texttt{(}\ \opt{\nt{params}}\ \texttt{)}\ \texttt{:}\
                     \nt{block}\\
\nt{params}  &\bnfeq \nt{param}\ \rep{\texttt{,}\ \nt{param}}\\
\nt{param}   &\bnfeq \nt{ident}\ \opt{\texttt{:}\ \nt{type}}\
                     \opt{\texttt{=}\ \nt{expr}}\\[4pt]
\nt{block}   &\bnfeq \texttt{\{}\ \rep{\nt{decl}}\ \texttt{\}}\\[4pt]
\nt{importDecl}&\bnfeq \tm{import}\ \nt{ident}\ \opt{\tm{as}\ \nt{ident}}\\
             &\bnfor \tm{from}\ \nt{ident}\ \tm{import}\
                     \nt{ident}\ \rep{\texttt{,}\ \nt{ident}}\\[4pt]
\nt{projectDecl}&\bnfeq \tm{project}\ \nt{ident}\ \nt{block}
\end{align*}

\subsection{The scientific-reasoning layer}

\begin{align*}
\nt{propDecl}  &\bnfeq \tm{proposition}\ \nt{ident}\ \texttt{:}\
                       \nt{propBody}\\
\nt{propBody}  &\bnfeq \texttt{\{}\ \rep{\nt{motionDecl}}\ \rep{\nt{stmt}}\
                       \texttt{\}}\\
\nt{motionDecl}&\bnfeq \tm{motion}\ \nt{ident}\
                       \texttt{(}\ \nt{string}\ \texttt{)}\\[4pt]
\nt{supportStmt}&\bnfeq \tm{support}\ \nt{ident}\
                       \opt{\tm{with\_confidence}\ \texttt{(}\ \nt{expr}\
                       \texttt{)}}\\
\nt{contradictStmt}&\bnfeq \tm{contradict}\ \nt{ident}\
                       \opt{\tm{with\_confidence}\ \texttt{(}\ \nt{expr}\
                       \texttt{)}}\\
\nt{inconclusiveStmt}&\bnfeq \tm{inconclusive}\
                       \texttt{(}\ \nt{string}\ \texttt{)}\\[4pt]
\nt{evidenceDecl}&\bnfeq \tm{evidence}\ \nt{ident}\ \texttt{:}\
                       \tm{sources}\ \texttt{:}\ \rep{\nt{source}}\\
\nt{source}    &\bnfeq \texttt{-}\ \nt{expr}\\[4pt]
\nt{patternDecl}&\bnfeq \tm{pattern}\ \nt{ident}\ \texttt{:}\
                       \nt{patternBody}\\
\nt{patternBody}&\bnfeq \texttt{\{}\ \nt{signature}\ \opt{\nt{withinStmt}}\
                       \texttt{\}}\\
\nt{signature} &\bnfeq \texttt{signature}\ \texttt{:}\ \nt{string}\\[4pt]
\nt{goalDecl}  &\bnfeq \tm{goal}\ \nt{ident}\ \texttt{:}\ \nt{goalBody}\\
\nt{metaDecl}  &\bnfeq \tm{metacognitive}\ \nt{ident}\ \texttt{:}\
                       \nt{metaBody}
\end{align*}

\subsection{The probabilistic layer}

\begin{align*}
\nt{pointDecl} &\bnfeq \tm{point}\ \nt{ident}\ \texttt{=}\ \nt{pointLit}\\
\nt{pointLit}  &\bnfeq \texttt{\{}\ \nt{field}\
                       \rep{\texttt{,}\ \nt{field}}\ \texttt{\}}\\
\nt{field}     &\bnfeq \nt{ident}\ \texttt{:}\ \nt{expr}\\[4pt]
\nt{resolveStmt}&\bnfeq \tm{resolution}\ \nt{ident}\
                       \texttt{(}\ \opt{\nt{args}}\ \texttt{)}\\
               &\bnfor \tm{resolve}\ \nt{expr}\ \opt{\nt{ctx}}\\
\nt{ctx}       &\bnfeq \tm{given}\ \tm{context}\ \texttt{(}\ \nt{expr}\
                       \texttt{)}\\[4pt]
\nt{cycleStmt} &\bnfeq \tm{cycle}\ \nt{ident}\ \tm{over}\ \nt{expr}\
                       \texttt{:}\ \nt{block}\\
\nt{driftStmt} &\bnfeq \tm{drift}\ \nt{ident}\ \tm{in}\ \nt{expr}\
                       \texttt{:}\ \nt{block}\\
\nt{flowStmt}  &\bnfeq \tm{flow}\ \nt{ident}\ \tm{on}\ \nt{expr}\
                       \texttt{:}\ \nt{block}\\
\nt{rollStmt}  &\bnfeq \tm{roll}\ \tm{until}\ \tm{settled}\ \texttt{:}\
                       \nt{block}
\end{align*}

\subsection{Statements and control flow}

\begin{align*}
\nt{stmt}      &\bnfeq \nt{itemStmt} \bnfor \nt{assign} \bnfor \nt{givenStmt}
                       \bnfor \nt{withinStmt} \bnfor \nt{consideringStmt}\\
               &\bnfor \nt{forStmt} \bnfor \nt{whileStmt} \bnfor \nt{returnStmt}
                       \bnfor \nt{ensureStmt} \bnfor \nt{allowStmt}\\
               &\bnfor \nt{researchStmt} \bnfor \nt{causeStmt}
                       \bnfor \nt{supportStmt} \bnfor \nt{contradictStmt}\\
               &\bnfor \nt{inconclusiveStmt} \bnfor \nt{resolveStmt}
                       \bnfor \nt{cycleStmt} \bnfor \nt{driftStmt}\\
               &\bnfor \nt{flowStmt} \bnfor \nt{rollStmt} \bnfor \nt{tryStmt}
                       \bnfor \nt{exprStmt}\\[4pt]
\nt{itemStmt}  &\bnfeq \tm{item}\ \nt{ident}\ \opt{\texttt{:}\ \nt{type}}\
                       \texttt{=}\ \nt{expr}\\
\nt{assign}    &\bnfeq \nt{lvalue}\ \texttt{=}\ \nt{expr}\\
\nt{lvalue}    &\bnfeq \nt{ident} \bnfor \nt{expr}\,\texttt{.}\,\nt{ident}
                       \bnfor \nt{expr}\,\texttt{[}\,\nt{expr}\,\texttt{]}\\[4pt]
\nt{givenStmt} &\bnfeq \tm{given}\ \nt{expr}\ \texttt{:}\ \nt{block}\
                       \opt{\tm{given}\ \tm{otherwise}\ \texttt{:}\ \nt{block}}\\
\nt{withinStmt}&\bnfeq \tm{within}\ \nt{expr}\ \opt{\tm{as}\ \nt{ident}}\
                       \texttt{:}\ \nt{block}\\
\nt{consideringStmt}&\bnfeq \tm{considering}\ \nt{quant}\ \nt{ident}\ \tm{in}\
                       \nt{expr}\ \texttt{:}\ \nt{block}\\
\nt{quant}     &\bnfeq \texttt{all} \bnfor \texttt{these} \bnfor \texttt{item}
                       \bnfor \varepsilon\\
\nt{forStmt}   &\bnfeq \tm{for}\ \tm{each}\ \nt{ident}\ \tm{in}\ \nt{expr}\
                       \texttt{:}\ \nt{block}\\
\nt{whileStmt} &\bnfeq \tm{while}\ \nt{expr}\ \texttt{:}\ \nt{block}\\[4pt]
\nt{returnStmt}&\bnfeq \tm{return}\ \opt{\nt{expr}}\\
\nt{ensureStmt}&\bnfeq \tm{ensure}\ \nt{expr}\\
\nt{allowStmt} &\bnfeq \tm{allow}\ \nt{expr}\\
\nt{researchStmt}&\bnfeq \tm{research}\ \nt{expr}\\
\nt{causeStmt} &\bnfeq \tm{cause}\ \nt{ident}\ \texttt{=}\ \nt{expr}\\
\nt{tryStmt}   &\bnfeq \tm{try}\ \texttt{:}\ \nt{block}\
                       \rep{\nt{catch}}\ \opt{\tm{finally}\ \texttt{:}\
                       \nt{block}}\\
\nt{catch}     &\bnfeq \tm{catch}\ \opt{\nt{ident}}\ \opt{\tm{as}\ \nt{ident}}\
                       \texttt{:}\ \nt{block}\\
\nt{exprStmt}  &\bnfeq \nt{expr}
\end{align*}

\subsection{Expressions and precedence}
\label{sec:exprgrammar}

The expression grammar is stratified to encode operator precedence and
associativity, following the standard cascade \citep{aho2006}. Lower in the
cascade binds tighter. The pipe operators \texttt{|} and \texttt{|>} compose a
value with a function, threading the left operand as the first argument of the
right (cf.\ the functional-composition and dataflow traditions
\citep{wadler1992,backus1959}).

\begin{align*}
\nt{expr}       &\bnfeq \nt{pipe}\\
\nt{pipe}       &\bnfeq \nt{or}\ \rep{(\texttt{|}\bnfor\texttt{|>})\ \nt{or}}\\
\nt{or}         &\bnfeq \nt{and}\ \rep{\texttt{||}\ \nt{and}}\\
\nt{and}        &\bnfeq \nt{equality}\ \rep{\texttt{\&\&}\ \nt{equality}}\\
\nt{equality}   &\bnfeq \nt{comparison}\
                        \rep{(\texttt{==}\bnfor\texttt{!=})\ \nt{comparison}}\\
\nt{comparison} &\bnfeq \nt{term}\
                        \rep{(\texttt{<}\bnfor\texttt{>}\bnfor\texttt{<=}
                        \bnfor\texttt{>=})\ \nt{term}}\\
\nt{term}       &\bnfeq \nt{factor}\
                        \rep{(\texttt{+}\bnfor\texttt{-})\ \nt{factor}}\\
\nt{factor}     &\bnfeq \nt{unary}\
                        \rep{(\texttt{*}\bnfor\texttt{/})\ \nt{unary}}\\
\nt{unary}      &\bnfeq (\texttt{!}\bnfor\texttt{-})\ \nt{unary}
                        \bnfor \nt{call}\\
\nt{call}       &\bnfeq \nt{primary}\
                        \rep{\nt{callSuffix}}\\
\nt{callSuffix} &\bnfeq \texttt{(}\ \opt{\nt{args}}\ \texttt{)}
                        \bnfor \texttt{.}\ \nt{ident}
                        \bnfor \texttt{[}\ \nt{expr}\ \texttt{]}\\
\nt{args}       &\bnfeq \nt{expr}\ \rep{\texttt{,}\ \nt{expr}}\\
\nt{primary}    &\bnfeq \nt{ident}\bnfor\nt{num}\bnfor\nt{string}\bnfor\nt{bool}
                        \bnfor \texttt{(}\ \nt{expr}\ \texttt{)}\\
                &\bnfor \nt{arrayLit}\bnfor\nt{mapLit}\\
\nt{arrayLit}   &\bnfeq \texttt{[}\ \opt{\nt{args}}\ \texttt{]}\\
\nt{mapLit}     &\bnfeq \texttt{\{}\ \opt{\nt{field}\
                        \rep{\texttt{,}\ \nt{field}}}\ \texttt{\}}
\end{align*}
All binary operators are left-associative except assignment (\cref{sec:concrete}
treats assignment as a statement, not an expression, so no associativity
question arises). Unary operators are prefix and bind tighter than any binary
operator but looser than call/index/member suffixes.

\subsection{Types in the surface syntax}

\begin{align*}
\nt{type}  &\bnfeq \texttt{Str}\bnfor\texttt{Num}\bnfor\texttt{Bool}\bnfor
                   \texttt{Unit}\bnfor\texttt{Value}\\
           &\bnfor \texttt{TextUnit}\bnfor\texttt{Point}\bnfor\texttt{Points}
                   \bnfor\texttt{Entity}\bnfor\texttt{Floor}\bnfor\texttt{Corpus}\\
           &\bnfor \texttt{List}\,\texttt{[}\,\nt{type}\,\texttt{]}
                   \bnfor \texttt{Map}\,\texttt{[}\,\nt{type}\,\texttt{]}\\
           &\bnfor \nt{type}\ \texttt{=>}\ \nt{type}
\end{align*}

\begin{theorem}[Unambiguity of the core grammar]
\label{thm:unambiguous}
The grammar of \cref{sec:exprgrammar}, restricted to expressions, is
unambiguous, and the full statement grammar is unambiguous up to the
dangling-\tm{given} resolution that binds an \tm{otherwise} clause to the
nearest preceding \tm{given}.
\end{theorem}
\begin{proof}[Proof sketch]
The expression stratification assigns each operator a unique precedence level
and left-associativity, so every operator string has a unique parse by the
standard argument for precedence-cascade grammars \citep{aho2006}. Within a
level, the left-recursive repetition $\rep{op\ \nt{e}}$ forces left grouping.
The only statement-level ambiguity is the optional \tm{otherwise} clause, which
is resolved by the nearest-match (most-closely-nested) rule, exactly as for the
dangling-else of ALGOL-style languages \citep{naur1963}; with that rule fixed,
every sentential form has a unique derivation. The remaining productions are
$LL(1)$ after left-factoring of the \nt{stmt} alternatives, each of which is
selected by a distinct leading keyword.
\end{proof}

% ===========================================================================
\section{Abstract Syntax and Desugaring}
\label{sec:abstract}
% ===========================================================================

The concrete grammar contains redundancy for the programmer's convenience. We
reduce it to an \emph{abstract syntax} (a small core) by a structural
desugaring $\mathcal{D}$. Semantics in \cref{sec:dynamics} are defined on the
core.

\subsection{Core abstract syntax}

\begin{definition}[Core terms]
\label{def:core}
The core distinguishes \emph{expressions} $e$ and \emph{statements} $s$:
\begin{align*}
e &::= n \mid str \mid b \mid x \mid e \mathbin{\widehat{op}} e \mid
       \widehat{\neg}\, e \mid e(\overline{e}) \mid e.\ell \mid e[e]
       \mid [\overline{e}] \mid \{\overline{\ell : e}\}\\
  &\quad\mid \mathsf{point}\{\overline{\ell:e}\}
       \mid \mathsf{resolve}(e, \Strat) \mid \mathsf{textop}(\theta, e,
       \overline{e})\\[3pt]
s &::= \mathsf{item}\ x = e \mid x := e \mid \mathsf{seq}(\overline{s})
       \mid \mathsf{given}(e, s, s) \mid \mathsf{within}(e, x, s)\\
  &\quad\mid \mathsf{foreach}(x, e, s) \mid \mathsf{while}(e, s)
       \mid \mathsf{return}(e) \mid \mathsf{ensure}(e) \mid \mathsf{fun}(f,
       \overline{x{:}\tau}, s)\\
  &\quad\mid \mathsf{prop}(P, \overline{M{:}str}, s)
       \mid \mathsf{support}(M, e) \mid \mathsf{contradict}(M, e)
       \mid \mathsf{inconc}(str)\\
  &\quad\mid \mathsf{cycle}(x, e, s) \mid \mathsf{drift}(x, e, s)
       \mid \mathsf{flow}(x, e, s) \mid \mathsf{roll}(s)\\
  &\quad\mid \mathsf{consider}(\kappa, x, e, s)
       \mid \mathsf{import}(\overline{\cdot}) \mid \mathsf{try}(s,
       \overline{c}, s)
\end{align*}
where $x,f,P,M$ range over identifiers, $\ell$ over field labels, $\theta$ over
text-operation tags, $\Strat$ over resolution strategies (\cref{sec:resolution}),
$\kappa\in\{\mathsf{all},\mathsf{these},\mathsf{item}\}$ over the
\tm{considering} quantifiers, and $\widehat{op},\widehat{\neg}$ over the binary
and unary operators of \cref{sec:lexical}.
\end{definition}

\subsection{Desugaring}

The translation $\mathcal{D}$ from concrete to core is homomorphic on shared
constructs; the nontrivial clauses are:

\begin{itemize}[leftmargin=1.4em]
\item \textbf{Suites to blocks.} A colon-suite is mapped to a brace block
(\cref{rem:layout}); a sequence of declarations/statements becomes
$\mathsf{seq}(\dots)$.
\item \textbf{Conditionals.} $\mathsf{given}\ e: s$ without an \tm{otherwise}
desugars to $\mathsf{given}(e, \mathcal{D}\sem{s}, \mathsf{seq}())$; with an
\tm{otherwise} clause $s'$, to $\mathsf{given}(e, \mathcal{D}\sem{s},
\mathcal{D}\sem{s'})$.
\item \textbf{Considering quantifiers.} $\mathsf{considering}\ \kappa\ x\
\mathsf{in}\ e: s$ keeps $\kappa$; the default ($\varepsilon$) is
$\mathsf{item}$. The \tm{all}/\tm{these} forms differ semantically in whether
the body filters (\cref{sec:hybrid}).
\item \textbf{Pipes.} $e_1\ \texttt{|}\ e_2 \;\mapsto\; e_2(\mathcal{D}\sem{e_1})$
and $e_1\ \texttt{|>}\ e_2 \;\mapsto\; e_2(\mathcal{D}\sem{e_1})$ when $e_2$ is a
call target; if $e_2 = g(\overline{a})$ then the left operand is inserted as the
first argument: $g(\mathcal{D}\sem{e_1}, \overline{a})$.
\item \textbf{Function defaults.} A parameter $x:\tau = e_0$ desugars to a
guard at entry binding $x$ to $e_0$ when the corresponding argument is absent.
\item \textbf{Pattern and evidence declarations} desugar to a $\mathsf{prop}$
scaffold and a resource binding respectively (\cref{sec:propsem,sec:fourfile}).
\item \textbf{Text operators.} An application of a text-operation keyword (e.g.\
$\mathsf{divide}$, $\mathsf{summarize}$) on a \textsf{TextUnit} desugars to
$\mathsf{textop}(\theta, e, \overline{e})$.
\end{itemize}

\begin{proposition}[Desugaring is meaning-preserving by definition]
The dynamic semantics of \cref{sec:dynamics} is defined on core terms; the
meaning of a concrete program $p$ is by definition the meaning of
$\mathcal{D}\sem{p}$. Hence $\mathcal{D}$ is trivially adequate, and any two
concrete programs with equal desugaring are observationally equivalent.
\end{proposition}

% ===========================================================================
\section{Semantic Domains}
\label{sec:domains}
% ===========================================================================

We fix the mathematical objects over which programs compute. The development
follows the domain-theoretic tradition \citep{scott1976,stoy1977,winskel1993}
but stays elementary: all domains here are sets (with order where needed), and
recursion is handled operationally in \cref{sec:dynamics}.

\subsection{Values}
\label{sec:values}

\begin{definition}[Values]
\label{def:values}
The set of \emph{values} $\Val$ is the least set closed under the following
formation rules:
\begin{align*}
v ::=\ &\mathsf{str}(s) \mid \mathsf{num}(r) \mid \mathsf{bool}(b) \mid
        \mathsf{none}
      \mid \mathsf{list}(\overline{v}) \mid \mathsf{map}(m)\\
     &\mid \mathsf{clo}(\overline{x{:}\tau}, s, \env) \mid \mathsf{tu}(s, m)
      \mid \mathsf{pt}(p) \mid \mathsf{floor}(F) \mid \mathsf{ent}(p) \mid
      \mathsf{corpus}(s)
\end{align*}
where $s\in\mathit{String}$, $r\in\mathbb{R}$, $b\in\mathbb{B}=\{\top,\bot\}$,
$m\colon \mathit{Label}\rightharpoonup\Val$ a finite map, $\env$ an environment
(\cref{sec:envs}), $p$ a point (\cref{sec:points}), and $F$ a floor
(\cref{sec:points}). A \emph{text unit} $\mathsf{tu}(s,m)$ is a string with
metadata; an \emph{entity} $\mathsf{ent}(p)$ is a point that has been resolved
to determinate status.
\end{definition}

\subsection{The confidence algebra}
\label{sec:confalg}

Degrees of belief are elements of the unit interval. We equip $\Unit$ with two
binary operations that govern, respectively, the \emph{accumulation} of
independent corroborating evidence and the \emph{attenuation} of belief by
opposing evidence. These operations are the algebraic heart of the language's
treatment of uncertainty and are chosen to satisfy the axioms one expects of
such combination \citep{cox1946,shafer1976,jaynes2003}.

\begin{definition}[Confidence algebra]
\label{def:confalg}
The \emph{confidence algebra} is the structure
$\Conf = (\Unit, \nor, \att, 0, 1, \le)$ where $\le$ is the usual order and,
for $c,d\in\Unit$,
\[
c \nor d \;=\; 1-(1-c)(1-d), \jspace c \att d \;=\; c\,(1-d).
\]
The operation $\nor$ (\emph{noisy-or}, or probabilistic sum) accumulates
support; $\att$ attenuates a belief $c$ by a counter-belief $d$.
\end{definition}

\begin{proposition}[Algebraic laws]
\label{prop:conflaws}
For all $c,d,e\in\Unit$:
\begin{enumerate}[label=(\roman*),leftmargin=1.9em,nosep]
\item $(\Unit,\nor,0)$ is a commutative monoid: $\nor$ is associative,
commutative, and has unit $0$;
\item $\nor$ is idempotent-free but \emph{inflationary}: $c \le c\nor d$ and
$d \le c\nor d$, with $c\nor d \le 1$;
\item $\nor$ is monotone in each argument; $1$ is absorbing
($1\nor d = 1$);
\item $\att$ is antitone in its second argument and monotone in its first;
$c\att 0 = c$, $c \att 1 = 0$, and $0\le c\att d\le c$.
\end{enumerate}
\end{proposition}
\begin{proof}
(i) Writing $\bar c=1-c$, we have $c\nor d = 1-\bar c\bar d$, and
$\overline{c\nor d}=\bar c\bar d$; since multiplication on $[0,1]$ is
commutative, associative, with unit $1$, the map $c\mapsto\bar c$ is an
isomorphism from $(\Unit,\nor,0)$ to $([0,1],\cdot,1)$, transporting all three
laws. (ii) $c\nor d - c = (1-c)d \ge 0$, and $c\nor d\le 1$ because
$\bar c\bar d\ge 0$. (iii) $\partial(c\nor d)/\partial c = 1-d\ge 0$; and
$1\nor d = 1-0=1$. (iv) $\partial(c\att d)/\partial d=-c\le0$,
$\partial(c\att d)/\partial c = 1-d\ge0$; the boundary values are immediate.
\end{proof}

\begin{remark}[Why noisy-or]
Among associative, commutative, monotone combination rules on $\Unit$ with unit
$0$ and top $1$ that treat sources as independent, the noisy-or is the unique
one arising from the probability of at least one of independent Bernoulli
events firing \citep{pearl1988}. It is the combination rule of the language's
\emph{maximum-likelihood} aggregation strategy (\cref{sec:propsem}); other
strategies use other combinators, all definable from $\nor,\att$, and the
order.
\end{remark}

\subsection{Points, interpretations, floors, corpora}
\label{sec:points}

The probabilistic layer rests on the notion of a \emph{point}: a unit of
meaning carrying its own uncertainty and a distribution over readings. This is
the language-level reification of a fuzzy or ambiguous datum
\citep{zadeh1965,hajek1998}.

\begin{definition}[Interpretation, point]
\label{def:point}
An \emph{interpretation} is a pair $\iota=(m, p)$ with $m\in\mathit{String}$ a
candidate meaning and $p\in\Unit$ its probability. A \emph{point} is a tuple
\[
p \;=\; \big(\mathit{content}\in\mathit{String},\;
              \mathit{conf}\in\Unit,\;
              I\in \mathit{Interp}^{*},\;
              (\ell,u)\in\Unit^2 \big)
\]
where $I=\overline{(m_i,p_i)}$ is a finite list of interpretations subject to
the sub-stochastic constraint $\sum_i p_i \le 1$, and $(\ell,u)$ with
$\ell\le u$ are semantic bounds. The residual mass $1-\sum_i p_i$ is the
probability of an \emph{unlisted} reading.
\end{definition}

\begin{definition}[Interpretation entropy and ambiguity]
\label{def:entropy}
For a point $p$ with interpretations $I=\overline{(m_i,p_i)}$ and $k=|I|\ge1$,
the \emph{interpretation entropy} is the Shannon entropy
\citep{shannon1948} of the normalised distribution
$\hat p_i = p_i/\sum_j p_j$,
\[
H(p) \;=\; -\sum_{i=1}^{k}\hat p_i \log_2 \hat p_i,
\]
and the \emph{ambiguity} is $A(p)=H(p)/\log_2 k\in\Unit$ (with $A(p)=0$ when
$k=1$). A point's \emph{intrinsic confidence} satisfies
$\mathit{conf}(p) = 1 - A(p)$ when not otherwise specified.
\end{definition}

\begin{definition}[Floor, corpus]
\label{def:floor}
A \emph{floor} is a finite map $F\colon \mathit{Key}\rightharpoonup
(\mathit{Point}\times\mathbb{R}_{\ge0})$ assigning to each key a point and a
nonnegative weight; it is the language's \emph{probabilistic collection}, an
iterable dictionary of weighted points. A \emph{corpus} is a value
$\mathsf{corpus}(s)$ wrapping a body of text $s$ to be processed by stochastic
traversal. The \emph{weight-normalised} iteration order of $F$ visits each
$(k, p, w)$ with selection mass $w/\sum_{k'} w_{k'}$.
\end{definition}

\subsection{Environments, stores, and the trace}
\label{sec:envs}

\begin{definition}[Environment]
An \emph{environment} $\env\colon \mathit{Var}\rightharpoonup\Val$ is a finite
partial map from identifiers to values. The empty environment is $\varnothing$;
update is $\env[x\mapsto v]$; lookup is $\env(x)$.
\end{definition}

Two distinguished components accompany the environment during the evaluation of
a proposition or a hybrid loop and record, respectively, the state of an
ongoing argument and the auxiliary confidences produced by probabilistic
iteration.

\begin{definition}[Debate state]
\label{def:debate}
A \emph{debate state} is a finite map
$\Delta\colon \mathit{Motion}\rightharpoonup
(\mathcal{M}_{\mathrm{f}}(\mathit{Ev})\times\mathcal{M}_{\mathrm{f}}
(\mathit{Ev}))$ assigning to each motion a pair $(\mathsf{Aff},\mathsf{Con})$ of
finite multisets of \emph{evidence items}. An evidence item is a triple
$\eta=(\mathit{strength}\in\Unit, \mathit{conf}\in\Unit, \mathit{weight}\in
\mathbb{R}_{\ge0})$.
\end{definition}

\begin{definition}[Configuration]
A \emph{configuration} is a tuple $\langle \env, \Delta\rangle$. Evaluation
threads a configuration; ordinary statements leave $\Delta$ unchanged, while
$\mathsf{support}$/$\mathsf{contradict}$ statements extend it.
\end{definition}

% ===========================================================================
\section{The Type System}
\label{sec:types}
% ===========================================================================

Turbulance is dynamically typed with optional annotations and inference. We
formalise this as a \emph{gradual} type system in the sense of
\citet{siek2006}: a dynamic type $\dyn$ (written \texttt{Value} in source)
together with a \emph{consistency} relation $\consistent$ lets annotated and
unannotated code interoperate, while the fully annotated fragment enjoys the
soundness theorem of \cref{sec:meta}.

\subsection{Types and their order}

\begin{definition}[Types]
\label{def:types}
The set of types is generated by
\[
\tau ::= \mathsf{Str}\mid\mathsf{Num}\mid\mathsf{Bool}\mid\mathsf{Unit}
        \mid \mathsf{TextUnit} \mid \mathsf{Point}\mid\mathsf{Points}
        \mid\mathsf{Entity}\mid\mathsf{Floor}\mid\mathsf{Corpus}
        \mid \mathsf{List}[\tau]\mid\mathsf{Map}[\tau]\mid \tau\to\tau \mid
        \dyn.
\]
$\dyn$ is the dynamic (unknown) type. We call a type \emph{static} if $\dyn$
does not occur in it.
\end{definition}

\begin{definition}[Subtyping]
\label{def:subtype}
The subtype relation $\subtype$ is the least preorder containing:
\[
\mathsf{Entity}\subtype\mathsf{Point}, \jspace
\mathsf{Point}\subtype\mathsf{Points}\ \text{is \emph{not} assumed},\jspace
\mathsf{Points} \equiv \mathsf{List}[\mathsf{Point}],
\]
together with the covariant rules
$\dfrac{\tau\subtype\tau'}{\mathsf{List}[\tau]\subtype\mathsf{List}[\tau']}$,
$\dfrac{\tau\subtype\tau'}{\mathsf{Map}[\tau]\subtype\mathsf{Map}[\tau']}$, and
the standard contravariant/covariant rule for functions
$\dfrac{\tau_1'\subtype\tau_1 \quad \tau_2\subtype\tau_2'}
       {\tau_1\to\tau_2 \subtype \tau_1'\to\tau_2'}$.
A resolved point is a point (hence $\mathsf{Entity}\subtype\mathsf{Point}$); a
\textsf{Points} value is a homogeneous list of points.
\end{definition}

\begin{definition}[Consistency]
\label{def:consistency}
The \emph{consistency} relation $\consistent$ is the least symmetric relation
with
\[
\tau\consistent\tau,\jspace \dyn\consistent\tau,\jspace \tau\consistent\dyn,
\]
closed under the same structural rules as $\subtype$ (covariantly on
constructors, contra/co on $\to$). Consistency is \emph{not} transitive;
$\dyn\consistent\mathsf{Num}$ and $\dyn\consistent\mathsf{Str}$ do not give
$\mathsf{Num}\consistent\mathsf{Str}$.
\end{definition}

We write $\tau \lesssim \tau'$ for ``$\tau$ is consistent with a subtype of
$\tau'$'': $\exists\tau''.\ \tau\consistent\tau''\subtype\tau'$. This is the
relation used at function application and assignment.

\subsection{Typing judgements: expressions}

A typing context $\Gam\colon\mathit{Var}\rightharpoonup\tau$ assigns types to
variables. The expression judgement is $\Gam\vdash e:\tau$.

\begin{gather*}
\infrule{T-Num}{}{\Gam\vdash n:\mathsf{Num}} \jspace
\infrule{T-Str}{}{\Gam\vdash str:\mathsf{Str}} \jspace
\infrule{T-Bool}{}{\Gam\vdash b:\mathsf{Bool}} \\[10pt]
\infrule{T-Var}{x\in\dom(\Gam)}{\Gam\vdash x:\Gam(x)} \jspace
\infrule{T-None}{}{\Gam\vdash \mathsf{none}:\mathsf{Unit}}\\[10pt]
\infrule{T-Arith}{\Gam\vdash e_1:\tau_1 \quad \tau_1\lesssim\mathsf{Num}\quad
                  \Gam\vdash e_2:\tau_2 \quad \tau_2\lesssim\mathsf{Num}\quad
                  op\in\{+,-,*,/\}}
                 {\Gam\vdash e_1\mathbin{op}e_2:\mathsf{Num}}\\[10pt]
\infrule{T-Cmp}{\Gam\vdash e_1:\tau_1\quad\Gam\vdash e_2:\tau_2\quad
                \tau_1\consistent\tau_2\quad op\in\{<,>,\le,\ge,==,!=\}}
               {\Gam\vdash e_1\mathbin{op}e_2:\mathsf{Bool}}\\[10pt]
\infrule{T-Logic}{\Gam\vdash e_i:\tau_i\quad\tau_i\lesssim\mathsf{Bool}\
                  (i=1,2)\quad op\in\{\&\&,||\}}
                 {\Gam\vdash e_1\mathbin{op}e_2:\mathsf{Bool}}\\[10pt]
\infrule{T-App}{\Gam\vdash e:\tau_f\quad \tau_f\consistent
                (\tau_1\!\to\!\tau_2)\quad \Gam\vdash e_i:\tau_i'\quad
                \tau_i'\lesssim\tau_1\ \text{(pointwise)}}
               {\Gam\vdash e(\overline{e_i}):\tau_2}\\[10pt]
\infrule{T-Idx}{\Gam\vdash e:\tau\quad \tau\lesssim\mathsf{List}[\tau']\quad
                \Gam\vdash e':\tau''\quad\tau''\lesssim\mathsf{Num}}
               {\Gam\vdash e[e']:\tau'} \jspace
\infrule{T-Mem}{\Gam\vdash e:\tau\quad \tau\lesssim\mathsf{Map}[\tau']}
              {\Gam\vdash e.\ell:\tau'}\\[10pt]
\infrule{T-Pipe}{\Gam\vdash e_1:\tau_1\quad\Gam\vdash e_2:\tau_2\quad
                 \tau_2\consistent(\tau_1\!\to\!\tau_3)}
                {\Gam\vdash e_1\mathbin{|}e_2:\tau_3}\\[10pt]
\infrule{T-Point}{\Gam\vdash e_{\mathit{content}}:\tau_c\quad
                  \tau_c\lesssim\mathsf{Str}\quad
                  \Gam\vdash e_{\mathit{conf}}:\tau_k\quad
                  \tau_k\lesssim\mathsf{Num}}
                 {\Gam\vdash \mathsf{point}\{\dots\}:\mathsf{Point}}\\[10pt]
\infrule{T-Resolve}{\Gam\vdash e:\tau\quad\tau\lesssim\mathsf{Point}}
                   {\Gam\vdash \mathsf{resolve}(e,\Strat):\mathsf{Entity}}
\end{gather*}
The rule \rn{T-Resolve} captures the central typing fact of the probabilistic
layer: \emph{resolving a point yields an entity}---a value whose ambiguity has
been discharged.

\subsection{Typing judgements: statements}

The statement judgement is $\Gam\vdash s\dashv\Gam'$: under $\Gam$, statement
$s$ is well-typed and \emph{produces} the extended context $\Gam'$ (binding
introductions flow rightward).

\begin{gather*}
\infrule{T-Item}{\Gam\vdash e:\tau'\quad \tau'\lesssim\tau}
               {\Gam\vdash (\mathsf{item}\ x{:}\tau = e)\dashv\Gam[x\mapsto\tau]}
\jspace
\infrule{T-ItemDyn}{\Gam\vdash e:\tau}
               {\Gam\vdash (\mathsf{item}\ x = e)\dashv\Gam[x\mapsto\tau]}\\[10pt]
\infrule{T-Assign}{x\in\dom(\Gam)\quad \Gam\vdash e:\tau'\quad
                   \tau'\lesssim\Gam(x)}
                  {\Gam\vdash (x := e)\dashv\Gam}\\[10pt]
\infrule{T-Seq}{\Gam\vdash s_1\dashv\Gam_1 \quad
                \Gam_1\vdash s_2\dashv\Gam_2}
               {\Gam\vdash \mathsf{seq}(s_1,s_2)\dashv\Gam_2}\\[10pt]
\infrule{T-Given}{\Gam\vdash e:\tau\quad\tau\lesssim\mathsf{Bool}\quad
                  \Gam\vdash s_1\dashv\Gam_1\quad\Gam\vdash s_2\dashv\Gam_2}
                 {\Gam\vdash\mathsf{given}(e,s_1,s_2)\dashv\Gam}\\[10pt]
\infrule{T-Foreach}{\Gam\vdash e:\tau\quad\tau\lesssim\mathsf{List}[\tau']\quad
                    \Gam[x\mapsto\tau']\vdash s\dashv\Gam'}
                   {\Gam\vdash\mathsf{foreach}(x,e,s)\dashv\Gam}\\[10pt]
\infrule{T-While}{\Gam\vdash e:\tau\quad\tau\lesssim\mathsf{Bool}\quad
                  \Gam\vdash s\dashv\Gam'}
                 {\Gam\vdash\mathsf{while}(e,s)\dashv\Gam}\\[10pt]
\infrule{T-Within}{\Gam\vdash e:\tau\quad
                   \Gam[x\mapsto\tau]\vdash s\dashv\Gam'}
                  {\Gam\vdash\mathsf{within}(e,x,s)\dashv\Gam}\\[10pt]
\infrule{T-Ensure}{\Gam\vdash e:\tau\quad\tau\lesssim\mathsf{Bool}}
                  {\Gam\vdash\mathsf{ensure}(e)\dashv\Gam}\jspace
\infrule{T-Return}{\Gam\vdash e:\tau}
                  {\Gam\vdash\mathsf{return}(e)\dashv\Gam}\\[10pt]
\infrule{T-Fun}{\Gam[\overline{x\mapsto\tau}]\vdash s\dashv\Gam'\quad
                \Gam'\ \text{returns}\ \tau_r}
               {\Gam\vdash\mathsf{fun}(f,\overline{x{:}\tau},s)\dashv
                \Gam[f\mapsto\overline{\tau}\to\tau_r]}
\end{gather*}
The judgement ``$\Gam'$ returns $\tau_r$'' holds when every $\mathsf{return}(e)$
on a control path through the body has $\Gam'\vdash e:\tau_e$ with
$\tau_e\lesssim\tau_r$, and $\tau_r$ is the least upper bound (in $\subtype$) of
those $\tau_e$, or $\mathsf{Unit}$ if there is no return.

\subsection{Typing the scientific-reasoning layer}

\begin{gather*}
\infrule{T-Prop}{\overline{M}\ \text{distinct}\quad
                 \Gam,\ \overline{M{:}\mathsf{Motion}}\vdash s\dashv\Gam'}
                {\Gam\vdash\mathsf{prop}(P,\overline{M{:}str},s)\dashv
                 \Gam[P\mapsto\mathsf{Prop}]}\\[10pt]
\infrule{T-Support}{\Gam(M)=\mathsf{Motion}\quad
                    \Gam\vdash e:\tau\quad\tau\lesssim\mathsf{Num}}
                   {\Gam\vdash\mathsf{support}(M,e)\dashv\Gam}\jspace
\infrule{T-Contra}{\Gam(M)=\mathsf{Motion}\quad
                    \Gam\vdash e:\tau\quad\tau\lesssim\mathsf{Num}}
                   {\Gam\vdash\mathsf{contradict}(M,e)\dashv\Gam}
\end{gather*}
A $\mathsf{support}$ or $\mathsf{contradict}$ statement is well-typed only
inside the scope of a proposition that declares the named motion; the
type \textsf{Motion} is in scope exactly within \rn{T-Prop}'s premise. The
confidence expression, when present, must be numeric (and is clamped to
$\Unit$ at runtime; \cref{sec:propsem}).

\subsection{Typing the hybrid iteration constructs}

\begin{gather*}
\infrule{T-Cycle}{\Gam\vdash e:\tau\quad\tau\lesssim\mathsf{Floor}\quad
                  \Gam[x\mapsto\mathsf{Point}]\vdash s\dashv\Gam'}
                 {\Gam\vdash\mathsf{cycle}(x,e,s)\dashv\Gam}\jspace
\infrule{T-Drift}{\Gam\vdash e:\tau\quad\tau\lesssim\mathsf{Corpus}\quad
                  \Gam[x\mapsto\mathsf{Str}]\vdash s\dashv\Gam'}
                 {\Gam\vdash\mathsf{drift}(x,e,s)\dashv\Gam}\\[10pt]
\infrule{T-Flow}{\Gam\vdash e:\tau\quad\tau\lesssim\mathsf{Floor}\quad
                 \Gam[x\mapsto\mathsf{Str}]\vdash s\dashv\Gam'}
                {\Gam\vdash\mathsf{flow}(x,e,s)\dashv\Gam}\jspace
\infrule{T-Roll}{\Gam\vdash s\dashv\Gam'}
               {\Gam\vdash\mathsf{roll}(s)\dashv\Gam}\\[10pt]
\infrule{T-Consider}{\Gam\vdash e:\tau\quad\tau\lesssim\mathsf{Corpus}\
                     \text{or}\ \tau\lesssim\mathsf{List}[\tau'']\quad
                     \Gam[x\mapsto\tau_x]\vdash s\dashv\Gam'}
                    {\Gam\vdash\mathsf{consider}(\kappa,x,e,s)\dashv\Gam}
\end{gather*}
where $\tau_x=\mathsf{Str}$ when $\tau\lesssim\mathsf{Corpus}$ and
$\tau_x=\tau''$ otherwise. Inside a hybrid loop the implementation binds the
distinguished auxiliary variables \texttt{current\_weight},
\texttt{current\_confidence}, \texttt{current\_mode}, \texttt{current\_settled}
(\cref{sec:hybrid}); these are typed \textsf{Num}, \textsf{Num}, \textsf{Str},
\textsf{Bool} respectively in $\Gam'$ within $s$.

% ===========================================================================
\section{Operational Semantics}
\label{sec:dynamics}
% ===========================================================================

We give a big-step (natural) operational semantics \citep{kahn1987,winskel1993}.
Expression evaluation has the form
$\langle e,\env\rangle \evalto \langle v,\env'\rangle$, allowing calls to have
effects on the environment; statement evaluation has the form
$\langle s,\env,\Delta\rangle \evalto \langle\kappa,\env',\Delta'\rangle$,
where the \emph{completion} $\kappa\in\{\nrm\}\cup\{\ret\,v : v\in\Val\}$
records whether the statement finished normally or returned. A configuration
that matches no rule \emph{goes wrong}, written $\wrong$; \cref{sec:meta} shows
the static fragment never goes wrong.

\subsection{Expressions}
\label{sec:exprsem}

\begin{gather*}
\infrule{E-Lit}{}{\langle \mathit{lit},\env\rangle\evalto\langle
                  \sem{\mathit{lit}},\env\rangle}\jspace
\infrule{E-Var}{x\in\dom(\env)}{\langle x,\env\rangle\evalto\langle
                  \env(x),\env\rangle}\\[10pt]
\infrule{E-Bin}{\langle e_1,\env\rangle\evalto\langle v_1,\env_1\rangle\quad
                \langle e_2,\env_1\rangle\evalto\langle v_2,\env_2\rangle\quad
                v=\sem{op}(v_1,v_2)}
               {\langle e_1\mathbin{op}e_2,\env\rangle\evalto\langle
                v,\env_2\rangle}\\[10pt]
\infrule{E-App}{\langle e,\env\rangle\evalto\langle
                \mathsf{clo}(\overline{x{:}\tau},s,\env_0),\env_1\rangle\quad
                \langle \overline{e_i},\env_1\rangle\evalto\langle
                \overline{v_i},\env_2\rangle}
               {\langle e(\overline{e_i}),\env\rangle\evalto\langle
                v,\env_2\rangle}\\
\text{where}\quad
\langle s,\ \env_0[\overline{x\mapsto v_i}],\ \Delta_\bot\rangle\evalto
\langle \ret\,v,\ \env_3,\ \Delta_\bot\rangle,\quad
\text{or}\ v=\mathsf{none}\ \text{if completion is }\nrm.\\[10pt]
\infrule{E-Idx}{\langle e,\env\rangle\evalto\langle\mathsf{list}(\overline{v}),
                \env_1\rangle\quad
                \langle e',\env_1\rangle\evalto\langle\mathsf{num}(i),
                \env_2\rangle\quad 0\le i<|\overline{v}|}
               {\langle e[e'],\env\rangle\evalto\langle v_{\lfloor i\rfloor},
                \env_2\rangle}
\end{gather*}
Here $\sem{op}$ denotes the strict mathematical operation on values (numeric
arithmetic, comparison returning $\mathsf{bool}$, logical connectives with the
usual short-circuit reading for \texttt{\&\&},\texttt{||}); $\Delta_\bot$ is the
empty debate state used when no proposition is in scope. \rn{E-App} uses
call-by-value and lexical scope: the closure's captured environment $\env_0$ is
extended with the arguments.

\subsection{Statements and control flow}

\begin{gather*}
\infrule{S-Item}{\langle e,\env\rangle\evalto\langle v,\env_1\rangle}
               {\langle \mathsf{item}\ x=e,\env,\Delta\rangle\evalto
                \langle \nrm,\env_1[x\mapsto v],\Delta\rangle}\\[10pt]
\infrule{S-Seq1}{\langle s_1,\env,\Delta\rangle\evalto\langle\nrm,\env_1,
                 \Delta_1\rangle\quad
                 \langle s_2,\env_1,\Delta_1\rangle\evalto\langle\kappa,\env_2,
                 \Delta_2\rangle}
                {\langle \mathsf{seq}(s_1,s_2),\env,\Delta\rangle\evalto
                 \langle\kappa,\env_2,\Delta_2\rangle}\\[10pt]
\infrule{S-Seq2}{\langle s_1,\env,\Delta\rangle\evalto\langle\ret\,v,\env_1,
                 \Delta_1\rangle}
                {\langle \mathsf{seq}(s_1,s_2),\env,\Delta\rangle\evalto
                 \langle\ret\,v,\env_1,\Delta_1\rangle}\\[10pt]
\infrule{S-GivenT}{\langle e,\env\rangle\evalto\langle\mathsf{bool}(\top),
                   \env_1\rangle\quad
                   \langle s_1,\env_1,\Delta\rangle\evalto\langle\kappa,\env_2,
                   \Delta_2\rangle}
                  {\langle\mathsf{given}(e,s_1,s_2),\env,\Delta\rangle\evalto
                   \langle\kappa,\env_2,\Delta_2\rangle}\\[10pt]
\infrule{S-GivenF}{\langle e,\env\rangle\evalto\langle\mathsf{bool}(\bot),
                   \env_1\rangle\quad
                   \langle s_2,\env_1,\Delta\rangle\evalto\langle\kappa,\env_2,
                   \Delta_2\rangle}
                  {\langle\mathsf{given}(e,s_1,s_2),\env,\Delta\rangle\evalto
                   \langle\kappa,\env_2,\Delta_2\rangle}\\[10pt]
\infrule{S-Ret}{\langle e,\env\rangle\evalto\langle v,\env_1\rangle}
              {\langle\mathsf{return}(e),\env,\Delta\rangle\evalto\langle
               \ret\,v,\env_1,\Delta\rangle}\\[10pt]
\infrule{S-EnsT}{\langle e,\env\rangle\evalto\langle\mathsf{bool}(\top),
                 \env_1\rangle}
                {\langle\mathsf{ensure}(e),\env,\Delta\rangle\evalto\langle
                 \nrm,\env_1,\Delta\rangle}\jspace
\infrule{S-EnsF}{\langle e,\env\rangle\evalto\langle\mathsf{bool}(\bot),
                 \env_1\rangle}
                {\langle\mathsf{ensure}(e),\env,\Delta\rangle\evalto\wrong}
\end{gather*}
The rule \rn{S-EnsF} gives $\mathsf{ensure}$ its meaning as an assertion in the
sense of \citet{hoare1969}: a failed ensure aborts the computation. The
$\mathsf{given}$ construct is the language's guarded command
\citep{dijkstra1975}; with the desugaring of \cref{sec:abstract} a missing
\tm{otherwise} branch is the no-op $\mathsf{seq}()$.

\paragraph{Bounded and unbounded iteration.}
For $\mathsf{foreach}(x,e,s)$ with $\langle e,\env\rangle\evalto\langle
\mathsf{list}(v_1\cdots v_n),\env_1\rangle$, evaluation runs $s$ once per
element, threading the configuration and stopping early on a $\ret$:
\[
\infrule{S-For}{\langle e,\env\rangle\evalto\langle\mathsf{list}(\overline{v}),
                \env_1\rangle\quad
                \langle s, \env_1[x\mapsto v_1],\Delta\rangle\evalto\cdots
                \evalto\langle\kappa,\env_n,\Delta_n\rangle}
               {\langle\mathsf{foreach}(x,e,s),\env,\Delta\rangle\evalto
                \langle\kappa',\env_n,\Delta_n\rangle}
\]
where $\kappa'$ is $\ret\,v$ if any iteration returned (the first such, after
which iteration stops) and $\nrm$ otherwise. The rule for $\mathsf{while}$ is
the usual unfolding and need not terminate.

\subsection{Within: scoped structural traversal}

The $\mathsf{within}(e,x,s)$ construct evaluates $e$ to a structured value and
runs $s$ with $x$ bound to it, providing a scope for pattern-guarded actions;
when $e$ denotes a collection the body is run once per constituent. It is the
structural counterpart of $\mathsf{foreach}$ specialised to semantic units
(text units, points): the body typically contains $\mathsf{given}$ guards that
fire $\mathsf{support}$/$\mathsf{contradict}$.

\[
\infrule{S-Within}{\langle e,\env\rangle\evalto\langle v,\env_1\rangle\quad
                   \langle s,\env_1[x\mapsto v],\Delta\rangle\evalto
                   \langle\kappa,\env_2,\Delta_2\rangle}
                  {\langle\mathsf{within}(e,x,s),\env,\Delta\rangle\evalto
                   \langle\kappa,\env_2,\Delta_2\rangle}
\]

\subsection{Propositions: an argumentation semantics}
\label{sec:propsem}

Evaluating a proposition runs its body in a fresh debate state, then scores each
motion under a resolution strategy and aggregates to a proposition verdict. This
realises inference as a defeasible argument \citep{toulmin1958,dung1995,
besnard2008} rather than a deduction.

\begin{gather*}
\infrule{S-Prop}{\Delta_0=[\,\overline{M\mapsto(\varnothing,\varnothing)}\,]\quad
                 \langle s,\env,\Delta_0\rangle\evalto\langle\_,\env_1,
                 \Delta_1\rangle}
                {\langle\mathsf{prop}(P,\overline{M{:}str},s),\env,\Delta\rangle
                 \evalto\langle\nrm,\env_1[P\mapsto
                 \mathsf{verdict}_\Strat(\Delta_1)],\Delta\rangle}
\end{gather*}
The body's $\mathsf{support}$ and $\mathsf{contradict}$ statements extend the
\emph{inner} debate state $\Delta_0$:
\begin{gather*}
\infrule{S-Sup}{\langle e,\env\rangle\evalto\langle\mathsf{num}(c),\env_1
                \rangle\quad c'=\mathrm{clip}_{\Unit}(c)\quad
                \eta=(c',c',1)}
               {\langle\mathsf{support}(M,e),\env,\Delta\rangle\evalto
                \langle\nrm,\env_1,\Delta[M\mathbin{+\!\!=}_{\mathsf{Aff}}\eta]
                \rangle}\\[10pt]
\infrule{S-Con}{\langle e,\env\rangle\evalto\langle\mathsf{num}(c),\env_1
                \rangle\quad c'=\mathrm{clip}_{\Unit}(c)\quad
                \eta=(c',c',1)}
               {\langle\mathsf{contradict}(M,e),\env,\Delta\rangle\evalto
                \langle\nrm,\env_1,\Delta[M\mathbin{+\!\!=}_{\mathsf{Con}}\eta]
                \rangle}
\end{gather*}
where $\Delta[M\mathbin{+\!\!=}_{\mathsf{Aff}}\eta]$ adds $\eta$ to the
affirmation multiset of $M$ (symmetrically for $\mathsf{Con}$), and the default
confidence (when $\tm{with\_confidence}$ is omitted) is $c'=1$.

\begin{definition}[Motion score and proposition verdict]
\label{def:score}
Fix a strategy $\Strat$. For a motion $M$ with
$\Delta(M)=(\mathsf{Aff},\mathsf{Con})$ define the \emph{motion score}
$\mathrm{score}_\Strat(M)\in\Unit$. For the four basic strategies, writing
$a_i=(\mathit{str}_i,\mathit{cf}_i,w_i)$ for affirmations and $b_j$ for
contentions, with $W=\sum_i w_i+\sum_j w_j$:
\begin{align*}
\textsc{MaxLik}:\quad& \mathrm{score}=\mathrm{clip}_\Unit\!\Big(
   \tfrac{1}{W}\big(\textstyle\sum_i w_i\,\mathit{str}_i\,\mathit{cf}_i
   -\sum_j w_j\,\mathit{str}_j\,\mathit{cf}_j\big)\Big),\\
\textsc{Bayes}:\quad& \mathrm{score}=\textstyle\prod_i
   \mathrm{post}(\,\cdot\,;a_i)\ \text{starting from prior } \tfrac12,\
   \text{updated by } a_i,b_j\ \text{as likelihoods},\\
\textsc{Conserv}:\quad& \mathrm{score}=\max\!\big(0,\ \textsc{MaxLik}
   -\lambda\textstyle\sum_j w_j\,\mathit{str}_j\big),\\
\textsc{Explor}:\quad& \mathrm{score}=\min\!\big(1,\ \textsc{MaxLik}
   +\mu\textstyle\sum_i w_i\,\mathit{str}_i\big),
\end{align*}
with fixed penalty/boost constants $\lambda,\mu\in\Unit$. The \emph{verdict} of
a motion is
\[
\mathrm{verdict}(M)=\begin{cases}
\textsf{Supported} & \mathrm{score}_\Strat(M)\ge\theta_+\\
\textsf{Contradicted} & \mathrm{score}_\Strat(M)\le\theta_-\\
\textsf{Inconclusive} & \text{otherwise,}
\end{cases}
\]
for thresholds $0\le\theta_-<\theta_+\le1$, and the \emph{proposition verdict}
$\mathrm{verdict}_\Strat(\Delta)$ aggregates the motions: it is the pair
$\big(\overline{(M,\mathrm{score}_\Strat(M),\mathrm{verdict}(M))},\,
\bar{c}\big)$ with overall confidence $\bar c=\mathrm{score}$ combined across
motions by the strategy's combinator ($\nor$ for \textsc{MaxLik}).
\end{definition}

\begin{remark}[Faithfulness to a debate]
Definition~\ref{def:score} is a quantitative instance of Dung's abstract
argumentation \citep{dung1995}: affirmations and contentions are arguments,
$\mathsf{contradict}$ is an attack, and the score is a graded acceptability
under the chosen semantics. The strategy is an explicit parameter, making the
basis of every conclusion inspectable.
\end{remark}

\subsection{Points and resolution}
\label{sec:resolution}

A \emph{resolution strategy} is a function that collapses a point's
interpretation distribution to a determinate outcome
\citep{kozen1981,mciver2005}. We give the canonical strategies; each maps a
point $p$ with normalised interpretations $\overline{(m_i,\hat p_i)}$ to a
\emph{resolution result}.

\begin{definition}[Resolution results and strategies]
\label{def:resolution}
A \emph{resolution result} is one of
\[
\mathsf{Certain}(v)\ \mid\ \mathsf{Uncertain}(\overline{(v_i,\hat p_i)},
(\ell,u),\bar c)\ \mid\ \mathsf{Fuzzy}(\phi,\bar t,\sigma),
\]
where $\phi$ is a membership function $\Val\to\Unit$, $\bar t$ its central
tendency, $\sigma$ its spread. For a fixed acceptance threshold $\theta$:
\begin{align*}
\textsc{MaxLik}(p)&=\begin{cases}
  \mathsf{Certain}(m_{i^\ast}) & \hat p_{i^\ast}\ge\theta,\ i^\ast=\arg\max_i
   \hat p_i,\\
  \mathsf{Uncertain}(\overline{(m_i,\hat p_i)},\mathrm{CI},1-A(p)) &
   \text{otherwise;}\end{cases}\\
\textsc{Bayes}(p)&=\textsc{MaxLik}\ \text{on the posterior } \hat p_i\propto
   \hat p_i\cdot\mathrm{prior}(m_i);\\
\textsc{Conserv}(p)&=\mathsf{Certain}(m_{i^\circ}),\quad
   i^\circ=\arg\min_i \mathrm{risk}(m_i)\ \text{(safest reading)};\\
\textsc{Explor}(p)&=\mathsf{Certain}(m_{i^\bullet}),\quad
   i^\bullet=\arg\max_i \mathrm{info}(m_i)\ \text{(most informative)};\\
\textsc{Full}(p)&=\mathsf{Uncertain}(\overline{(m_i,\hat p_i)},\mathrm{CI},
   1-A(p)).
\end{align*}
\end{definition}

\begin{gather*}
\infrule{S-Resolve}{\langle e,\env\rangle\evalto\langle\mathsf{pt}(p),\env_1
                    \rangle\quad
                    \mathsf{res}=\Strat(p)\quad v=\mathsf{ent}(\widehat{p})}
                   {\langle\mathsf{resolve}(e,\Strat),\env,\Delta\rangle
                    \evalto\langle\nrm,\env_1[\_\mapsto v],\Delta\rangle}
\end{gather*}
where $\widehat p$ is $p$ with its interpretation set collapsed according to
$\mathsf{res}$ (a singleton for $\mathsf{Certain}$, the full set otherwise) and
its confidence updated to the result's $\bar c$. Resolving an already-determinate
entity is the identity.

\subsection{The four hybrid iteration constructs}
\label{sec:hybrid}

The four constructs $\mathsf{cycle}$, $\mathsf{drift}$, $\mathsf{flow}$,
$\mathsf{roll}$ are instances of one schema: bind a loop variable and a set of
auxiliary confidences, run the body, accumulate. They differ in the
\emph{source} they traverse and in which auxiliaries they expose. We write
$\env\!\Downarrow$ to abbreviate threading through the body.

\paragraph{Cycle (deterministic traversal of a floor).}
$\mathsf{cycle}(x,e,s)$ evaluates $e$ to a floor $F$ and runs $s$ once per
entry in weight-normalised order, binding $x$ to the point's content and
\texttt{current\_weight} to the entry weight.
\[
\infrule{S-Cycle}{\langle e,\env\rangle\evalto\langle\mathsf{floor}(F),\env_1
                  \rangle\quad
                  \overline{(k_t,p_t,w_t)}=\mathrm{iter}(F)\quad
                  \langle s,\env_1[x\mapsto\mathsf{str}(\mathit{content}(p_t)),
                  \texttt{current\_weight}\mapsto\mathsf{num}(w_t)],\Delta\rangle
                  \evalto\cdots}
                 {\langle\mathsf{cycle}(x,e,s),\env,\Delta\rangle\evalto
                  \langle\nrm,\env_n,\Delta_n\rangle}
\]

\paragraph{Drift (stochastic traversal of a corpus).}
$\mathsf{drift}(x,e,s)$ evaluates $e$ to a corpus, draws a sequence of text
fragments by the implementation's sampling procedure, and runs $s$ per fragment,
binding $x$ to the fragment, \texttt{current\_confidence} to the fragment's
running confidence, and \texttt{current\_mode} to the processing mode. Sampling
is the only source of nondeterminism in the language; we model it as an oracle
$\mathcal{O}$ that the judgement is parameterised by, so that
$\langle s,\dots\rangle\evalto_\mathcal{O}\langle\dots\rangle$ is deterministic
given $\mathcal{O}$ (cf.\ \cref{thm:determinacy}).
\[
\infrule{S-Drift}{\langle e,\env\rangle\evalto\langle\mathsf{corpus}(s_0),\env_1
                  \rangle\quad \overline{f_t}=\mathcal{O}(s_0)\quad
                  \langle s,\env_1[x\mapsto\mathsf{str}(f_t),
                  \texttt{current\_confidence}\mapsto\mathsf{num}(c_t)],\Delta
                  \rangle\evalto\cdots}
                 {\langle\mathsf{drift}(x,e,s),\env,\Delta\rangle\evalto_
                  \mathcal{O}\langle\nrm,\env_n,\Delta_n\rangle}
\]

\paragraph{Flow (streaming over the lines of a source).}
$\mathsf{flow}(x,e,s)$ treats its source as a sequence of lines, running $s$ per
line and binding \texttt{current\_settled} to whether the line's processing has
converged. It is deterministic.

\paragraph{Roll (iterate to a fixpoint).}
$\mathsf{roll}(s)$ repeats $s$ until a convergence predicate holds, binding
\texttt{settled}, \texttt{final\_confidence}, and \texttt{iterations}. Let
$\Phi$ be the per-iteration update on the ambient confidence
$c\in\Unit$ induced by $s$ (a $\mathsf{resolve}/\mathsf{guess}$ step), and let
the convergence predicate be $\mathrm{settled}(c)\equiv c\ge\theta_{\mathrm{conv}}$
or a small-change criterion $|c'-c|<\epsilon$.
\[
\infrule{S-Roll}{c_0=\env(\texttt{current\_confidence})\quad
                 c_{t+1}=\Phi(c_t)\quad
                 N=\min\{t : \mathrm{settled}(c_t)\}}
                {\langle\mathsf{roll}(s),\env,\Delta\rangle\evalto
                 \langle\nrm,\env[\texttt{settled}\mapsto\mathsf{bool}(\top),
                 \texttt{final\_confidence}\mapsto\mathsf{num}(c_N),
                 \texttt{iterations}\mapsto\mathsf{num}(N)],\Delta\rangle}
\]
If no such $N$ exists the loop diverges (\cref{sec:meta} gives a sufficient
condition for $N$ to be finite).

\paragraph{Considering (probabilistic guarded iteration).}
$\mathsf{consider}(\kappa,x,e,s)$ iterates over the constituents of $e$ (the
sentences of a corpus, or the elements of a list), binds $x$, and runs $s$
under a per-item guard; with $\kappa=\mathsf{these}$ the body itself filters
(running only when its internal $\mathsf{given}$ succeeds), with
$\kappa=\mathsf{all}$ it runs on every item, and with $\kappa=\mathsf{item}$ on
a single distinguished item. The guard may be a \emph{confidence condition}
$\mathrm{conf}\mathbin{\bowtie}\theta$ with
$\bowtie\in\{>,<,\mathrm{within},\mathrm{outside}\}$, evaluated against the
running resolution confidence.

% ===========================================================================
\section{The Four-File Project Model}
\label{sec:fourfile}
% ===========================================================================

A Turbulance \emph{project} is not a single source file but a linked tuple of
four roles. We formalise the roles as typed modules and the project as their
composition under a linkage relation; this is a module system in the sense of
\citet{cardelli1985} specialised to four fixed roles.

\begin{definition}[Project]
\label{def:project}
A \emph{project} is a tuple
$\Pi=(\mathcal{T},\mathcal{G},\mathcal{F},\mathcal{H})$ where:
\begin{itemize}[leftmargin=1.4em,nosep]
\item $\mathcal{T}$ (the \emph{orchestration unit}, file role \texttt{.trb}) is
a core program (\cref{def:core}); it holds the functions, propositions, and the
control flow that drives the computation.
\item $\mathcal{G}$ (the \emph{dependency graph}, file role \texttt{.ghd}) is a
finite labelled directed graph $(V,E,\rho)$ whose vertices are named resources
and whose labelling $\rho$ binds each resource to an external source
(a dataset, a service, a knowledge base); it is the project's import surface.
\item $\mathcal{F}$ (the \emph{state surface}, file role \texttt{.fs}) is a
typed record of observable state and its presentation: a map from state keys to
$\Val$ together with a layout descriptor; it is the read-model the computation
publishes.
\item $\mathcal{H}$ (the \emph{decision log}, file role \texttt{.hre}) is an
append-only sequence of \emph{decision records}; a decision record is a tuple
$(\textit{time},\textit{decision},\textit{reasoning},\textit{confidence}\in\Unit,
\textit{evidence-refs})$. It is the project's audit trail.
\end{itemize}
\end{definition}

\begin{definition}[Linkage]
\label{def:linkage}
The four roles are linked by a relation $\ell\subseteq
\dom(\mathcal{T})\times(\dom(\mathcal{G})\cup\dom(\mathcal{F})\cup
\dom(\mathcal{H}))$ that resolves: (i) each free resource identifier of
$\mathcal{T}$ to a vertex of $\mathcal{G}$; (ii) each state-publishing operation
of $\mathcal{T}$ to a key of $\mathcal{F}$; and (iii) each proposition verdict
and resolution of $\mathcal{T}$ to an appended record of $\mathcal{H}$. A
project is \emph{well-linked} if $\ell$ is total on the free identifiers and
publishing/logging sites of $\mathcal{T}$.
\end{definition}

\begin{definition}[System state]
\label{def:sysstate}
During execution the project maintains a \emph{system state}
$\Sigma=(\env,\Delta,\mathcal{H}^{+},\mathcal{F}^{+})$ extending the
configuration of \cref{sec:envs} with the current decision log
$\mathcal{H}^{+}$ and published state $\mathcal{F}^{+}$. Evaluation of
$\mathcal{T}$ proceeds by the rules of \cref{sec:dynamics} with two additional
effects: a proposition verdict (rule \rn{S-Prop}) appends a record to
$\mathcal{H}^{+}$, and an explicit publish updates $\mathcal{F}^{+}$.
\end{definition}

\begin{proposition}[Separation of effects]
\label{prop:separation}
For a well-linked project, the orchestration semantics is conservative over the
core semantics of \cref{sec:dynamics}: erasing $\mathcal{H}^{+}$ and
$\mathcal{F}^{+}$ from every configuration yields exactly the core evaluation of
$\mathcal{T}$. Hence the decision log and state surface are \emph{observational}
---they record but do not alter the computation.
\end{proposition}
\begin{proof}
The two additional effects (append to $\mathcal{H}^{+}$, update
$\mathcal{F}^{+}$) are write-only with respect to the core: no core rule reads
$\mathcal{H}^{+}$ or $\mathcal{F}^{+}$. Therefore the projection that drops
these components commutes with every rule, and the diagram of core rules is
recovered exactly.
\end{proof}

% ===========================================================================
\section{Metatheory}
\label{sec:meta}
% ===========================================================================

We establish the basic properties a language standard must guarantee:
determinacy, type safety on the static fragment, soundness of the evidence
semantics, and termination criteria.

\subsection{Determinacy}

\begin{theorem}[Determinacy]
\label{thm:determinacy}
The deterministic fragment of Turbulance (all constructs except $\mathsf{drift}$
and $\mathsf{consider}$ over a corpus) has a deterministic big-step semantics:
if $\langle s,\env,\Delta\rangle\evalto\langle\kappa_1,\env_1,\Delta_1\rangle$
and $\langle s,\env,\Delta\rangle\evalto\langle\kappa_2,\env_2,\Delta_2\rangle$
then $(\kappa_1,\env_1,\Delta_1)=(\kappa_2,\env_2,\Delta_2)$. The full language
is deterministic relative to a fixed sampling oracle $\mathcal{O}$.
\end{theorem}
\begin{proof}
By induction on the height of the first derivation. For each syntactic form
exactly one rule applies, and that rule's premises are evaluations of
strictly smaller subterms (for expressions) or of the same statement on a
strictly smaller remaining structure (for the bounded loops), so the inductive
hypothesis gives uniqueness of each premise's result; the conclusion is a
function of those results. The only branching is between value-discriminating
rule pairs (\rn{S-GivenT}/\rn{S-GivenF}, \rn{S-EnsT}/\rn{S-EnsF},
\rn{S-Seq1}/\rn{S-Seq2}), and these are mutually exclusive because a value is
exactly one of $\mathsf{bool}(\top)$, $\mathsf{bool}(\bot)$, etc., and a
completion is exactly one of $\nrm$, $\ret\,v$. The constructs $\mathsf{drift}$
and corpus-$\mathsf{consider}$ consult $\mathcal{O}$; fixing $\mathcal{O}$
removes the only choice, and the same induction applies.
\end{proof}

\subsection{Type safety}

We prove safety for the \emph{static fragment}: programs all of whose
annotations are static types (no $\dyn$) and in which every position that the
rules require to be consistent is in fact related by $\subtype$. The argument is
the syntactic one of \citet{wright1994,pierce2002}. We first relate runtime
values to types.

\begin{definition}[Value typing]
$\vdash v:\tau$ holds when: $\vdash\mathsf{num}(r):\mathsf{Num}$;
$\vdash\mathsf{str}(s):\mathsf{Str}$; $\vdash\mathsf{bool}(b):\mathsf{Bool}$;
$\vdash\mathsf{none}:\mathsf{Unit}$; $\vdash\mathsf{list}(\overline v):
\mathsf{List}[\tau]$ if each $\vdash v_i:\tau$; $\vdash\mathsf{pt}(p):
\mathsf{Point}$; $\vdash\mathsf{ent}(p):\mathsf{Entity}$;
$\vdash\mathsf{floor}(F):\mathsf{Floor}$; $\vdash\mathsf{corpus}(s):
\mathsf{Corpus}$; $\vdash\mathsf{clo}(\overline{x{:}\tau},s,\env):
\overline\tau\to\tau_r$ if the closure body checks at $\tau_r$ under a context
typing $\env$. An environment is well-typed, $\vdash\env:\Gam$, if
$\vdash\env(x):\Gam(x)$ for all $x\in\dom(\Gam)$.
\end{definition}

\begin{lemma}[Canonical forms]
\label{lem:canonical}
If $\vdash v:\mathsf{Num}$ then $v=\mathsf{num}(r)$; if
$\vdash v:\mathsf{Bool}$ then $v=\mathsf{bool}(b)$; if
$\vdash v:\mathsf{List}[\tau]$ then $v=\mathsf{list}(\overline w)$ with each
$\vdash w_i:\tau$; if $\vdash v:\tau_1\to\tau_2$ then $v$ is a closure; if
$\vdash v:\mathsf{Point}$ then $v\in\{\mathsf{pt}(p),\mathsf{ent}(p)\}$; if
$\vdash v:\mathsf{Floor}$ then $v=\mathsf{floor}(F)$; if
$\vdash v:\mathsf{Corpus}$ then $v=\mathsf{corpus}(s)$.
\end{lemma}
\begin{proof}
Immediate from the definition of value typing, which assigns each base/structured
type a unique value former, and from $\mathsf{Entity}\subtype\mathsf{Point}$.
\end{proof}

\begin{theorem}[Preservation]
\label{thm:preservation}
In the static fragment, if $\Gam\vdash e:\tau$ and $\vdash\env:\Gam$ and
$\langle e,\env\rangle\evalto\langle v,\env'\rangle$, then $\vdash v:\tau$ and
$\vdash\env':\Gam$. Likewise if $\Gam\vdash s\dashv\Gam'$, $\vdash\env:\Gam$,
and $\langle s,\env,\Delta\rangle\evalto\langle\kappa,\env'',\Delta''\rangle$,
then $\vdash\env'':\Gam'$, and if $\kappa=\ret\,v$ then $\vdash v:\tau_r$ for the
body's return type.
\end{theorem}
\begin{proof}[Proof sketch]
By induction on the evaluation derivation, inverting the typing derivation at
each step. The base cases (\rn{E-Lit}, \rn{E-Var}) are immediate, the latter
from $\vdash\env:\Gam$. For \rn{E-Bin} with an arithmetic operator, the typing
rule \rn{T-Arith} forces both operands to be $\mathsf{Num}$ in the static
fragment; the IH gives $\vdash v_i:\mathsf{Num}$, hence by
\cref{lem:canonical} both are $\mathsf{num}(\cdot)$, $\sem{op}$ is defined and
returns a $\mathsf{num}$, of type $\mathsf{Num}$. For \rn{E-App}, \rn{T-App} in
the static fragment gives $\tau_f=\overline{\tau_1}\to\tau_2$ and arguments at
$\subtype\,\overline{\tau_1}$; the IH and \cref{lem:canonical} make $v$ a
closure whose body checks at $\tau_2$; evaluating the body under the extended
well-typed environment preserves $\tau_2$ by the statement IH. The statement
cases use \rn{T-Item}/\rn{S-Item} (the produced environment extends $\Gam$
coherently), \rn{T-Given}/\rn{S-Given$\cdot$} (both branches produce $\Gam$),
and \rn{T-Seq}/\rn{S-Seq$\cdot$} (threading). The scientific and probabilistic
rules preserve types because \rn{S-Prop} writes a $\mathsf{Prop}$ to $P$,
\rn{S-Resolve} writes an $\mathsf{Entity}$, and \rn{S-Sup}/\rn{S-Con} only
extend $\Delta$, leaving $\Gam$ fixed.
\end{proof}

\begin{theorem}[Progress]
\label{thm:progress}
In the static fragment, if $\Gam\vdash e:\tau$ and $\vdash\env:\Gam$ then either
$e$ is a value or there is a (unique, by \cref{thm:determinacy}) configuration
$\langle v,\env'\rangle$ with $\langle e,\env\rangle\evalto\langle v,\env'
\rangle$; in particular $e$ does not go wrong. The analogous statement holds for
statements, whose only non-value, non-progressing outcome would be a failed
$\mathsf{ensure}$ (rule \rn{S-EnsF})---which is a \emph{defined} abort, not a
stuck state.
\end{theorem}
\begin{proof}[Proof sketch]
The potential failures are: applying a non-function, indexing a non-list,
arithmetic on a non-number, or a missing variable. Each is excluded in the
static fragment by the corresponding typing rule together with
\cref{lem:canonical}: e.g.\ \rn{T-App} gives the operator a function type, so by
canonical forms its value is a closure and \rn{E-App} applies; \rn{T-Var}
requires $x\in\dom(\Gam)$ and $\vdash\env:\Gam$ gives $x\in\dom(\env)$. The only
remaining non-normal outcome is \rn{S-EnsF}, which is a deliberate, defined
abort in the sense of \citet{hoare1969}, not an undefined ``wrong''.
\end{proof}

\begin{corollary}[Safety of the static fragment]
A well-typed program in the static fragment never goes wrong: it either
produces a value/normal completion, deliberately aborts via a failed
$\mathsf{ensure}$, or diverges.
\end{corollary}

\begin{remark}[The dynamic fragment]
At positions typed $\dyn$, consistency replaces subtyping and the guarantees of
\cref{thm:preservation,thm:progress} are recovered \emph{up to a runtime type
error}: an implementation inserts a check at each $\dyn$-to-static boundary, and
a violated check raises a defined error rather than going wrong. This is the
gradual guarantee of \citet{siek2006}: fully annotated code is sound, fully
unannotated code runs with dynamic checks, and the two interoperate.
\end{remark}

\subsection{Soundness of the evidence semantics}

\begin{theorem}[Boundedness and monotonicity]
\label{thm:evidence}
For every strategy $\Strat\in\{\textsc{MaxLik},\textsc{Conserv},\textsc{Explor},
\textsc{Bayes}\}$ and motion $M$:
\begin{enumerate}[label=(\roman*),leftmargin=1.9em,nosep]
\item $\mathrm{score}_\Strat(M)\in\Unit$ (boundedness);
\item adding an affirmation $\eta$ to $\mathsf{Aff}(M)$ does not decrease
$\mathrm{score}_{\textsc{MaxLik}}(M)$, and adding a contention does not increase
it (monotonicity);
\item if $\mathsf{Con}(M)=\varnothing$ then $\mathrm{score}_{\textsc{MaxLik}}(M)
=\bigoplus_i (\mathit{str}_i\mathit{cf}_i)$ up to weight normalisation, so
corroborating evidence accumulates by the noisy-or of \cref{def:confalg}.
\end{enumerate}
\end{theorem}
\begin{proof}
(i) Each strategy clips into $\Unit$ (Def.~\ref{def:score}); \textsc{Bayes}
returns a posterior probability, already in $\Unit$. (ii) For \textsc{MaxLik},
$\mathrm{score}=\mathrm{clip}_\Unit(\frac1W(S_+-S_-))$ with
$S_+=\sum_i w_i\mathit{str}_i\mathit{cf}_i$ and $S_-$ the analogous contention
sum. Adding an affirmation increases $S_+$ and $W$; writing $f(S_+,W)=
\mathrm{clip}(\frac{S_+-S_-}{W})$, and noting each new term contributes
$w\,\mathit{str}\,\mathit{cf}\le w$, the normalised score is nondecreasing
because the added mass is weighted at least as favourably as the average; the
contention case is symmetric. (iii) With no contentions the unclipped score is
the weighted mean of $\mathit{str}_i\mathit{cf}_i$; under the noisy-or reading
of independent corroboration (Remark after Def.~\ref{def:confalg}) this is the
$\nor$-aggregate, which is the unique independent-source combinator
(Prop.~\ref{prop:conflaws}).
\end{proof}

\begin{proposition}[Confidence never exceeds certainty]
\label{prop:noovercert}
No sequence of $\mathsf{support}$ statements drives a motion's
\textsc{MaxLik} score above $1$ or a $\mathsf{resolve}$ result's confidence
above $1$. Belief is conserved within $\Unit$.
\end{proposition}
\begin{proof}
Immediate from Prop.~\ref{prop:conflaws}(ii)--(iii) ($\nor$ is bounded by $1$
and $1$ is absorbing) and the clipping in Def.~\ref{def:score} and
Def.~\ref{def:resolution}.
\end{proof}

\subsection{Termination}

\begin{theorem}[Termination of bounded iteration]
\label{thm:termbounded}
Evaluation of $\mathsf{foreach}$, $\mathsf{cycle}$, $\mathsf{flow}$, and
$\mathsf{consider}$ over a finite source terminates whenever the body
terminates on each element. $\mathsf{while}$ and $\mathsf{roll}$ may diverge.
\end{theorem}
\begin{proof}
The first four iterate over a finite list/floor/line-set of known length $n$;
the derivation is a sequence of $n$ body evaluations, each finite by hypothesis,
so the whole is finite. $\mathsf{while}$ is the standard unbounded loop, and
$\mathsf{roll}$ unfolds until a predicate holds (next theorem).
\end{proof}

\begin{theorem}[Convergence of \texorpdfstring{$\mathsf{roll}$}{roll}]
\label{thm:rollconv}
Let $\mathsf{roll}(s)$ have per-iteration confidence update $c_{t+1}=\Phi(c_t)$
with $\Phi\colon\Unit\to\Unit$ continuous and either (a) a contraction
($|\Phi(c)-\Phi(c')|\le L|c-c'|$ with $L<1$) toward a fixpoint $c^\ast\ge
\theta_{\mathrm{conv}}$, or (b) monotone nondecreasing with step bounded below
on $[c_0,\theta_{\mathrm{conv}})$. Then the convergence index
$N=\min\{t:\mathrm{settled}(c_t)\}$ is finite and \rn{S-Roll} terminates.
\end{theorem}
\begin{proof}
(a) By the Banach fixed-point theorem $c_t\to c^\ast$ geometrically, so
$|c_t-c^\ast|\le L^t|c_0-c^\ast|$; choosing $t$ with $L^t|c_0-c^\ast|<\epsilon$
makes the small-change criterion fire, so $N<\infty$. (b) If $\Phi$ is monotone
nondecreasing with $\Phi(c)-c\ge\delta>0$ for $c<\theta_{\mathrm{conv}}$, then
$c_t\ge c_0+t\delta$ exceeds $\theta_{\mathrm{conv}}$ after at most
$\lceil(\theta_{\mathrm{conv}}-c_0)/\delta\rceil$ steps, so $N$ is finite.
\end{proof}

\begin{remark}[Undecidability in general]
Without a hypothesis like \cref{thm:rollconv}, termination of $\mathsf{while}$
and $\mathsf{roll}$ is undecidable, as for any language with general recursion
\citep{hopcroft2006}; the language therefore provides bounded iteration
constructs ($\mathsf{foreach}$, $\mathsf{cycle}$) whose termination is
guaranteed, and reserves unbounded iteration for cases where the programmer
supplies a convergence argument.
\end{remark}

% ===========================================================================
\section{Worked Examples}
\label{sec:examples}
% ===========================================================================

\subsection{A proposition with motions and graded support}

The following program declares a hypothesis about a dataset, decomposes it into
three motions, and accumulates graded support by traversing the data with
$\mathsf{within}$ and $\mathsf{given}$ guards.

\begin{lstlisting}
proposition DataQuality:
    motion CompleteData("Dataset has no missing values")
    motion ConsistentFormat("All entries follow the schema")
    motion ReasonableValues("Values lie within expected ranges")

    within dataset:
        given missing_count == 0:
            support CompleteData with_confidence(0.99)
        given schema_valid:
            support ConsistentFormat with_confidence(0.95)
        given all_in_range(lo, hi):
            support ReasonableValues with_confidence(0.90)
        given outlier_rate > 0.2:
            contradict ReasonableValues with_confidence(0.6)
\end{lstlisting}

Evaluating \texttt{DataQuality} (rule \rn{S-Prop}) opens an inner debate state,
runs the $\mathsf{within}$ body once (rule \rn{S-Within}), and at each satisfied
guard appends an affirmation or contention (rules \rn{S-Sup}/\rn{S-Con}). Under
\textsc{MaxLik} with thresholds $\theta_+=0.7$, $\theta_-=0.3$ the motion
\texttt{ReasonableValues} receives one affirmation $(0.90,0.90,1)$ and possibly
one contention $(0.6,0.6,1)$; its score is
$\mathrm{clip}_\Unit\!\big(\tfrac{1}{2}(0.81-0.36)\big)=0.225$ if both fire,
yielding verdict \textsf{Inconclusive} (between $\theta_-$ and $\theta_+$), and
$0.81$ if only the affirmation fires, yielding \textsf{Supported}. The verdict,
its score, and the contributing evidence are appended to the decision log
$\mathcal{H}^{+}$ (\cref{def:sysstate}).

\subsection{A point and its resolution}

\begin{lstlisting}
point diagnosis = {
    content: "elevated marker with ambiguous aetiology",
    confidence: 0.7
}

item determined = resolve diagnosis given context(patient_history)
\end{lstlisting}

The literal builds $\mathsf{pt}(p)$ with $\mathit{conf}=0.7$ (rule
\rn{T-Point}); $\mathsf{resolve}$ applies the ambient strategy (rule
\rn{S-Resolve}), producing an $\mathsf{Entity}$ whose interpretation set is
collapsed---to a singleton if the dominant reading clears the acceptance
threshold $\theta$, otherwise to an $\mathsf{Uncertain}$ result carrying a
confidence interval and the residual ambiguity $1-A(p)$ (Def.~\ref{def:entropy},
Def.~\ref{def:resolution}). The type of \texttt{determined} is
$\mathsf{Entity}$ (rule \rn{T-Resolve}).

\subsection{A hybrid loop over a probabilistic floor}

\begin{lstlisting}
funxn assess(floor_of_claims):
    cycle claim over floor_of_claims:
        item r = resolve claim
        given current_weight > 0.5:
            support StrongClaims with_confidence(r.confidence)
    return StrongClaims
\end{lstlisting}

\texttt{cycle} (rule \rn{S-Cycle}) visits each weighted point of the floor in
weight-normalised order, binding \texttt{claim} and \texttt{current\_weight};
the body resolves the point and, when its selection weight is high, registers
graded support. Termination is guaranteed because a floor is finite
(\cref{thm:termbounded}).

\subsection{A four-file project, in outline}

A project \texttt{biomarker} comprises: \texttt{biomarker.trb} (the
orchestration unit $\mathcal{T}$ above); \texttt{biomarker.ghd} (the dependency
graph $\mathcal{G}$ binding \texttt{dataset}, \texttt{patient\_history}, and
external knowledge bases to sources); \texttt{biomarker.fs} (the state surface
$\mathcal{F}$ publishing running confidences and verdicts for display); and
\texttt{biomarker.hre} (the decision log $\mathcal{H}$ receiving each verdict and
resolution). By \cref{prop:separation} the computation's result is independent
of $\mathcal{F}$ and $\mathcal{H}$, which observe but do not steer it.

% ===========================================================================
\section{Related Paradigms}
\label{sec:related}
% ===========================================================================

Turbulance draws on, and can be located precisely against, several established
traditions.

\paragraph{Functional and imperative cores.} The expression language, closures,
and lexical scope are standard \citep{landin1966,church1940,barendregt1984};
the statement language with $\mathsf{item}$ bindings and structured control is
imperative in the tradition formalised by \citet{winskel1993}. The pipe
operators are the composition combinators of functional programming
\citep{wadler1992}.

\paragraph{Assertions and guarded commands.} $\mathsf{ensure}$ is a Hoare-style
assertion \citep{hoare1969}; $\mathsf{given}$ is a guarded command
\citep{dijkstra1975}, and the $\mathsf{within}/\mathsf{given}$ idiom is a
pattern-guarded structural recursion over semantic units.

\paragraph{Logic and argumentation.} Propositions and motions echo logic
programming's view of computation as proof search \citep{kowalski1974}, but the
semantics is \emph{graded} and \emph{defeasible}: it is a quantitative abstract
argumentation \citep{dung1995,besnard2008,toulmin1958} rather than a monotone
deduction, with $\mathsf{contradict}$ as attack and the resolution strategy as
the acceptability semantics.

\paragraph{Probabilistic and fuzzy computation.} Points, floors, and resolution
place Turbulance in the probabilistic-programming family \citep{kozen1981,
goodman2008,gordon2014,mciver2005}; the confidence algebra and the
membership-function results are fuzzy-set constructions \citep{zadeh1965,
hajek1998}; the accumulation of graded evidence is in the spirit of the
mathematical theory of evidence \citep{shafer1976,dempster1967} and Bayesian
confirmation \citep{cox1946,jaynes2003,pearl1988}, with interpretation entropy
measured \`a la \citet{shannon1948}.

\paragraph{Gradual typing.} The type system's dynamic type and consistency
relation are exactly those of \citet{siek2006}, instantiated with the language's
semantic types; subtyping follows \citet{cardelli1985}.

The novelty is not any single ingredient but their combination in one language
with a single control flow: a program may, in adjacent lines, bind a value,
accumulate graded support for a hypothesis, and resolve an ambiguous datum under
a named strategy---each with a precise, inspectable semantics.

% ===========================================================================
\section{Conclusion}
\label{sec:conclusion}
% ===========================================================================

We have specified Turbulance completely and from first principles: its lexical
and context-free grammars, its abstract syntax and desugaring, its semantic
domains (centred on a confidence algebra and the structures of points, floors,
and corpora), its gradual type system, and its big-step operational
semantics---including the argumentation semantics of propositions, the
resolution semantics of points, and the four hybrid iteration constructs. We
formalised the four-file project model as a typed module composition with a
linkage relation, and proved its decision log and state surface to be purely
observational. The metatheory establishes determinacy, type safety on the
static fragment, boundedness and monotonicity of the evidence semantics, and
termination criteria for the iteration constructs.

The specification is deliberately self-contained: every concept needed to
understand the language is defined here, and the only external references are to
the standard literature on which the design draws. It is intended to serve two
audiences. For the implementer it is a blueprint: the grammar fixes the parser,
the type rules fix the checker, and the operational rules fix the evaluator,
each with stated invariants. For the language user it is a standard: the meaning
of every construct is pinned down, and the guarantees---that beliefs stay within
$[0,1]$, that corroboration accumulates monotonically, that the static fragment
never goes wrong, that bounded loops terminate---are theorems, not conventions.

Two extensions are left to companion documents and do not affect the core
defined here: a denotational account that interprets the confidence algebra and
resolution strategies in a domain of sub-probability distributions, and an
operational treatment of the metacognitive and goal constructs as reflective
operations over the decision log. Neither is required to read, implement, or use
the language as specified above.

% ===========================================================================
\bibliographystyle{plainnat}
\bibliography{references}
% ===========================================================================

\end{document}